\section{AIME I 2023}
        \begin{enumerate}
        	\item (\textit{AIME I 2023}) Five men and nine women stand equally spaced around a circle in random order. The probability that every man stands diametrically opposite a woman is $\frac{m}{n},$ where $m$ and $n$ are relatively prime positive integers. Find $m+n.$


 
 \textbf{Answer: }191



\textbf{Solution 1}

For simplicity purposes, we consider two arrangements different even if they only differ by rotations or reflections. In this way, there are $14!$ arrangements without restrictions.

 First, there are $\binom{7}{5}$ ways to choose the man-woman diameters. Then, there are $10\cdot8\cdot6\cdot4\cdot2$ ways to place the five men each in a man-woman diameter. Finally, there are $9!$ ways to place the nine women without restrictions.

 Together, the requested probability is 
 \[\frac{\tbinom{7}{5}\cdot(10\cdot8\cdot6\cdot4\cdot2)\cdot9!}{14!} = \frac{21\cdot(10\cdot8\cdot6\cdot4\cdot2)}{14\cdot13\cdot12\cdot11\cdot10} = \frac{48}{143},\] from which the answer is $48+143 = \boxed{191}.$

 



\textbf{Solution 2 (Solution 4 but more efficient)}

Arranging the 14 people in a circle, one could just see that the probability depends on the men's placement, as any woman can be opposite to a man. Therefore, we just need to consider the men's arrangements.

 Place the first man on the circle. There are 14 ways to do this. Doing this removes two spots for the next man, the diametrically opposite slot and the slot already occupied, so the next man has 12 spots left to choose. The third man has 10, fourth man has 8, and fifth man has 6.

 In total, there are 14 slots for the first man, 13 for the second, 12 for the third, 11 for the fourth, and 10 for the fifth.

 We have \( P(nA) = \frac{14 \times 12 \times 10 \times 8 \times 6}{14 \times 13 \times 12 \times 11 \times 10} \Rightarrow \frac{8 \times 6}{13 \times 11} = \frac{48}{143} \).

 Our answer is simply \( 48 + 143 =\) $\boxed{191}.$

 



\textbf{Solution 3}

We can simply just loop through each of the men and find the probability that the person opposite from him is a woman.

 Start by sitting down the $1$st man. The probability that the person opposite to him is a woman is $\frac{9}{13}$ since out of the $13$ people who can sit opposite to him, $9$ can be a woman. With the $2$nd man, we can use the same logic: there are $11$ people who can sit opposite to him, but only $8$ of them are a woman, so the probability is $\frac{8}{11}.$ We use the same logic for the $3$rd, $4$th and $5$th men to get probabilities of $\frac{7}{9}$, $\frac{6}{7}$ and $\frac{5}{5},$ respectively.

 Multiplying these probabilities, we get a final answer of 
 \[\frac{9}{13}\cdot\frac{8}{11}\cdot\frac{7}{9}\cdot\frac{6}{7}\cdot\frac{5}{5}=\frac{48}{143}\longrightarrow 48+143 = \boxed{191}.\]

 



\textbf{Solution 4}

This problem is equivalent to solving for the probability that no man is standing diametrically opposite to another man. We can simply just construct this.

 We first place the $1$st man anywhere on the circle, now we have to place the $2$nd man somewhere around the circle such that he is not diametrically opposite to the first man. This can happen with a probability of $\frac{12}{13}$ because there are $13$ available spots, and $12$ of them are not opposite to the first man.

 We do the same thing for the $3$rd man, finding a spot for him such that he is not opposite to the other $2$ men, which would happen with a probability of $\frac{10}{12}$ using similar logic. Doing this for the $4$th and $5$th men, we get probabilities of $\frac{8}{11}$ and $\frac{6}{10}$ respectively.

 Multiplying these probabilities, we get, 
 \[\frac{12}{13}\cdot\frac{10}{12}\cdot\frac{8}{11}\cdot\frac{6}{10}=\frac{48}{143}\longrightarrow\boxed{191}.\]

 



\textbf{Solution 5}

Assume that rotations and reflections are distinct arrangements, and replace men and women with identical M's and W's, respectively. (We can do that because the number of ways to arrange $5$ men in a circle and the number of ways to arrange $9$ women in a circle, are constants.) The total number of ways to arrange $5$ M's and $9$ W's is $\binom{14}{5} = 2002.$ 

 To count the number of valid arrangements (i.e. arrangements where every M is diametrically opposite a W), we notice that exactly $2$ of the pairs of diametrically opposite positions must be occupied by $2$ W's. There are $\binom{7}{2} = 21$ ways to choose these $2$ pairs. For the remaining $5$ pairs, we have to choose which position is occupied by an M and which is occupied by a W. This can be done in $2^{5} = 32$ ways. Therefore, there are $21*32 = 672$ valid arrangements. 

 Therefore, the probability that an arrangement is valid is $\frac{672}{2002} = \frac{48}{143}$ for an answer of $\boxed{191}.$

 



\textbf{Solution 6}

To start off, we calculate the total amount of ways to organize all $14$ people irrespective of any constraints. This is simply ${14\choose5} = 2002$, because we just count how many ways we can place all $5$ men in any of the $14$ slots. 

 Since men cannot be diametrically opposite with each other, because of the constraints, placing down one man in any given spot will make another spot on the opposite side of the circle unable to hold any men. This means that placing down one man will effectively take away $2$ spots. 

 There are $14$ possible slots the first man can be placed. Once that man was placed, the next man only have $12$ possible slots because the slot that the first man is in is taken and the diametrically opposite spot to the first man can't have any men. Similar logic applies for the third man, who has $10$ possible slots. The fourth man has $8$ possible slots, and the fifth man has $6$ possible slots.

 This means the number of ways you can place all $5$ men down is $14 \cdot 12 \cdot 10 \cdot 8 \cdot 6$. However, since the men are all indistinct from each other, you also have to divide that value by $5! = 120$, since there are $120$ ways to arrange the $5$ men in each possible positioning of the men on the circle. This means the total number of ways to arrange the men around the circle so that none of them are diametrically opposite of each other is: $\frac{14 \cdot 12 \cdot 10 \cdot 8 \cdot 6}{5!} = 672$. The women simply fill in the rest of the available slots in each arrangement of men.

 Thus, the final probability is $\frac{672}{2002} = \frac{48}{143}$, meaning the answer is $48 + 143 = \boxed{191}$.

 



\textbf{Solution 7}

We will first assign seats to the men. The first man can be placed in any of the $14$ slots. The second man can be placed in any of the remaining $13$ seats, except for the one diametrically opposite to the first man. So, there are $13 - 1 = 12$ ways to seat him. With a similar argument, the third man can be seated in $10$ ways, the fourth man in $8$ ways and the last man in $6$ ways.  

 So, the total number of ways to arrange the men is $14 \cdot 12 \cdot 10 \cdot 8 \cdot 6$.    

 The women go to the remaining $9$ spots. Note that since none of the seats diametrically opposite to the men is occupied, each man is opposite a woman. The number of ways to arrange the women is therefore, simply $9!$, meaning that the total number of ways to arrange the people with restrictions is $14 \cdot 12 \cdot 10 \cdot 8 \cdot 6 \cdot 9!$ In general, there are $14!$ ways to arrange the people without restrictions. So, the probability is    
 \[\frac{14 \cdot 12 \cdot 10 \cdot 8 \cdot 6 \cdot 9!}{14!} = \frac{8 \cdot 6}{13 \cdot 11} = \frac{48}{143}.\]   The answer is $48 + 143 = \boxed{191}$.

 



\textbf{Solution 8}

First pin one man on one seat (to ensure no rotate situations). Then there are $13!$ arrangements.Because $5$ men must have women at their opposite side, we consider the $2$nd man and the woman opposite as one group and name it $P_2.$ There are $4$ groups, $P_1, P_2, P_3, P_4$ except the first man pinned on the same point. And for the rest $4$ women, name them $P_5$ and $P_6.$ First to order $P_1, P_2, P_3, P_4, P_5, P_6,$ there are $6!$ ways. For the $1$st man, there are $9$ women to choose, $8$ for the $2$nd, $\ldots,$ $5$ for the $5$th, and then for the $2$ women pairs $3$ and $1.$  Because every $2$ person in the group have chance to change their position, there are $2^6$ possibilities. 

 So the possibility is 
 \[P=\frac{6!\cdot 9\cdot 8\cdot 7\cdot 6\cdot 5\cdot 3\cdot 1 \cdot 2^6}{13!}=\frac{48}{143}.\]

 The answer is $48+143=\boxed{191}.$

 



\textbf{Solution 9}

We get around the condition that each man can't be opposite to another man by simply considering all $7$ diagonals, and choosing $5$ where there will be a single man. For each diagonal, the man can go on either side, and there are $\binom{14}{5}$ ways to arrange the men and the women in total. Thus our answer is $\frac{\binom{7}{5}\cdot 2^5}{\binom{14}{5}} = \frac{48}{143}.$We get $48 + 143 = \boxed{191}$

 



\textbf{Solution 10}

We can find the probability of one arrangement occurring, and multiply it by the total number of arrangements. 

 The probability of a man being in any specific position is $\frac{5}{14}.$ The probability of a woman being across from him is $\frac{9}{13}.$ The probability of a man being in any valid position is now $\frac{4}{12},$ and the probability of a woman being across from him is $\frac{8}{11},$ and so forth. We stop when there are no more men left. Multiplying these probabilities together, 

 
 \[P(\mathrm{One\ successful\ outcome})=\frac{5}{14}\cdot \frac{9}{13}\cdot \frac{4}{12}\cdot \frac{8}{11}\cdot \frac{3}{10}\cdot \frac{7}{9}\cdot \frac{2}{8}\cdot \frac{1}{6}\cdot \frac{6}{7}\cdot \frac{5}{5} = \frac{1}{2002}.\]To find the total number of successful outcomes, we consider the diagonals; the total number of diagonals to be made is $\binom75$, since there are $7$ total diagonals, and we want to choose $5$ of them to connect a man to a woman. For each of these diagonals, the man can be on either side of the diagonal.

 It follows that there are $2$ possibilities for each diagonal (man on one side, woman on the other, and vice versa). There are $5$ diagonals with a man and a woman, so there are $2^5$ different ways for these diagonals to appear. 

 There are $\binom75$ successful diagonals, and for each of these diagonals, there are $2^5$ ways to seat the men and the women, there are $\binom75$ $\cdot 2^5$ successful outcomes. 

 Recall that
 \[P(\mathrm{Any\ successful\ outcome})=P(\mathrm{One\ successful\ outcome})\cdot P(\mathrm{Total\ number\ of\ successful\ outcomes}).\]Therefore, 
 \[P(\mathrm{Any\ successful\ outcome}) = \frac{1}{2002}\cdot 2^5\cdot \binom75 = \frac{1}{2002}\cdot 2^5\cdot 21 = \frac{2^5\cdot 21}{2002} = \frac{48}{143}.\]The requested sum is $48+143=\boxed{191}.$

 -Benedict T (countmath1)

 


	\item (\textit{AIME I 2023}) Positive real numbers $b \not= 1$ and $n$ satisfy the equations 
 \[\sqrt{\log_b n} = \log_b \sqrt{n} \qquad \text{and} \qquad b \cdot \log_b n = \log_b (bn).\] The value of $n$ is $\frac{j}{k},$ where $j$ and $k$ are relatively prime positive integers. Find $j+k.$


 
 \textbf{Answer: }881



\textbf{Solution 1}

Denote $x = \log_b n$.Hence, the system of equations given in the problem can be rewritten as
 \begin{align*} \sqrt{x} & = \frac{1}{2} x , \\ bx & = 1 + x . \end{align*}
Solving the system gives $x = 4$ and $b = \frac{5}{4}$.Therefore,
 \[n = b^x = \frac{625}{256}.\]Therefore, the answer is $625 + 256 = \boxed{881}$.

 



\textbf{Solution 2}

We can use the property that $\log(xy) = \log(x) + \log(y)$ on the first equation. We get $\log_b(bn) = 1 + \log_b(n)$. Then, subtracting $\log_b(n)$ from both sides, we get $(b-1) \log_b(n) = 1$, therefore $\log_b(n) = \frac{1}{b-1}$. Substituting that into our first equation, we get $\frac{1}{2b-2} = \sqrt{\frac{1}{b-1}}$. Squaring, reciprocating, and simplifying both sides, we get the quadratic $4b^2 - 9b + 5 = 0$. Solving for $b$, we get $\frac{5}{4}$ and $1$. Since the problem said that $b \neq 1$, $b = \frac{5}{4}$. To solve for $n$, we can use the property that $\log_b(n) = \frac{1}{b-1}$. $\log_\frac{5}{4}(n) = 4$, so $n = \frac{5^4}{4^4} = \frac{625}{256}$. Adding these together, we get $\boxed{881}$

 



\textbf{Solution 3 (quick)}

There is some $x$ such that $n=b^{4x^2}$. The first equation becomes $2x=2x^2$, so $x=1$ since $x\ne 0$ following from the stated $b\ne 1$. Thus $n=b^{4\cdot 1^2}=b^4$, which turns the second equation into $4b=5$. Thus $b=\frac{5}{4}$ and $n= b^4=\frac{625}{256}$. Adding the numerator and denominator gives $\boxed{881}$.

 Honestly, this problem is kinda well placed.

 



\textbf{Solution 4}

Squaring the first equation, we have $\log_b{n}=\left(\dfrac{1}{2}\log_b{n}\right)^2=\dfrac{1}{4}(\log_b{n})^2$, so $\log_b{n}=4$ and $b^4=n$. The second equation can be rearranged as stated in earlier solutions to $(b-1)\log_b{n}=1$, so $n^{b-1}=b$. Substituting, we have $b^{4b-4}=b$, so $b=\dfrac{5}{4}$ and thus $n=\left(\dfrac{5}{4}\right)^4=\dfrac{625}{256}$, and thus the answer is $\boxed{881.}$

 



\textbf{Solution 5 (similar to solution 4)}

Start the same as solution 4. Square the first equation to get $\log_b{n}=\left(\dfrac{1}{2}\log_b{n}\right)^2$, and $\log_b{n}=4$. Then sub this into the second equation to get 4b = 4 + 1 = 5, so b =5/4. Since $\log_b{n}=4$, n = b\textasciicircum{}4, so n =$\dfrac{625}{256}$, and the answer is $\boxed{881.}$

 -Supernova77

 


	\item (\textit{AIME I 2023}) A plane contains $40$ lines, no $2$ of which are parallel. Suppose that there are $3$ points where exactly $3$ lines intersect, $4$ points where exactly $4$ lines intersect, $5$ points where exactly $5$ lines intersect, $6$ points where exactly $6$ lines intersect, and no points where more than $6$ lines intersect. Find the number of points where exactly $2$ lines intersect.


 
 \textbf{Answer: }607



\textbf{Solution}

In this solution, let \textbf{$\boldsymbol{n}$-line points} be the points where exactly $n$ lines intersect. We wish to find the number of $2$-line points.

 There are $\binom{40}{2}=780$ pairs of lines. Among them:

 
                \begin{itemize}
                    \item The $3$-line points account for $3\cdot\binom32=9$ pairs of lines.
                \end{itemize}


                
                \begin{itemize}
                    \item The $4$-line points account for $4\cdot\binom42=24$ pairs of lines.
                \end{itemize}


                
                \begin{itemize}
                    \item The $5$-line points account for $5\cdot\binom52=50$ pairs of lines.
                \end{itemize}


                
                \begin{itemize}
                    \item The $6$-line points account for $6\cdot\binom62=90$ pairs of lines.
                \end{itemize}


                It follows that the $2$-line points account for $780-9-24-50-90=\boxed{607}$ pairs of lines, where each pair intersect at a single point. 

 


	\item (\textit{AIME I 2023}) The sum of all positive integers $m$ such that $\frac{13!}{m}$ is a perfect square can be written as $2^a3^b5^c7^d11^e13^f,$ where $a,b,c,d,e,$ and $f$ are positive integers. Find $a+b+c+d+e+f.$


 
 \textbf{Answer: }012



\textbf{Solution 1}

We first rewrite $13!$ as a prime factorization, which is $2^{10}\cdot3^5\cdot5^2\cdot7\cdot11\cdot13.$

 For the fraction to be a square, it needs each prime to be an even power. This means $m$ must contain $7\cdot11\cdot13$. Also, $m$ can contain any even power of $2$ up to $2^{10}$, any odd power of $3$ up to $3^{5}$, and any even power of $5$ up to $5^{2}$. The sum of $m$ is 
 \[(2^0+2^2+2^4+2^6+2^8+2^{10})(3^1+3^3+3^5)(5^0+5^2)(7^1)(11^1)(13^1) =\] 
 \[1365\cdot273\cdot26\cdot7\cdot11\cdot13 = 2\cdot3^2\cdot5\cdot7^3\cdot11\cdot13^4.\] Therefore, the answer is $1+2+1+3+1+4=\boxed{012}$.

 



\textbf{Solution 2}

The prime factorization of $13!$ is 
 \[2^{10} \cdot 3^5 \cdot 5^2 \cdot 7 \cdot 11 \cdot 13.\]To get $\frac{13!}{m}$ a perfect square, we must have $m = 2^{2x} \cdot 3^{1 + 2y} \cdot 5^{2z} \cdot 7 \cdot 11 \cdot 13$, where $x \in \left\{ 0, 1, \cdots , 5 \right\}$, $y \in \left\{ 0, 1, 2 \right\}$, $z \in \left\{ 0, 1 \right\}$.

 Hence, the sum of all feasible $m$ is
 \begin{align*} \sum_{x=0}^5 \sum_{y=0}^2 \sum_{z=0}^1 2^{2x} \cdot 3^{1 + 2y} \cdot 5^{2z} \cdot 7 \cdot 11 \cdot 13 & = \left( \sum_{x=0}^5 2^{2x} \right) \left( \sum_{y=0}^2 3^{1 + 2y} \right) \left( \sum_{z=0}^1 5^{2z} \right) 7 \cdot 11 \cdot 13 \\ & = \frac{4^6 - 1}{4-1} \cdot \frac{3 \cdot \left( 9^3 - 1 \right)}{9 - 1} \cdot \frac{25^2  - 1}{25 - 1} \cdot 7 \cdot 11 \cdot 13 \\ &  = 2 \cdot 3^2 \cdot 5 \cdot 7^3 \cdot 11 \cdot 13^4 . \end{align*}


 Therefore, the answer is
 \begin{align*} 1 + 2 + 1 + 3 + 1 + 4 & = \boxed{012} . \end{align*}


 



\textbf{Solution 3 (Educated Guess and Engineer's Induction (Fake solve))}

Try smaller cases. There is clearly only one $m$ that makes $\frac{2!}{m}$ a square, and this is $m=2$. Here, the sum of the exponents in the prime factorization is just $1$. Furthermore, the only $m$ that makes $\frac{3!}{m}$ a square is $m = 6 = 2^13^1$, and the sum of the exponents is $2$ here. Trying $\frac{4!}{m}$ and $\frac{5!}{m}$, the sums of the exponents are $3$ and $4$. Based on this, we (incorrectly!) conclude that, when we are given $\frac{n!}{m}$, the desired sum is $n-1$. The problem gives us $\frac{13!}{m}$, so the answer is $13-1 = \boxed{012}$.

 -InsetIowa9

 However...

 The induction fails starting at $n = 9$ !The actual answers $f(n)$ for small $n$ are:$0, 1, 2, 3, 4, 5, 6, 7, 7, 7, 8, 11, 12$In general, $f(p) = f(p-1)+1$ if p is prime, $n=4,6,8$ are "lucky", and the pattern breaks down after $n=8$

 -"fake"  warning by oinava

 


 


	\item (\textit{AIME I 2023}) Let $P$ be a point on the circle circumscribing square $ABCD$ that satisfies $PA \cdot PC = 56$ and $PB \cdot PD = 90.$ Find the area of $ABCD.$


 
 \textbf{Answer: }106



\textbf{Solution 1 (Ptolemy's Theorem)}

\href{https://artofproblemsolving.com/wiki/index.php?title=Ptolemy\%27s\_theorem}{Ptolemy's theorem} states that for cyclic quadrilateral $WXYZ$, $WX\cdot YZ + XY\cdot WZ = WY\cdot XZ$.

 We may assume that $P$ is between $B$ and $C$. Let $PA = a$, $PB = b$, $PC = c$, $PD = d$, and $AB = s$. We have $a^2 + c^2 = AC^2 = 2s^2$, because $AC$ is a diameter of the circle. Similarly, $b^2 + d^2 = 2s^2$. Therefore, $(a+c)^2 = a^2 + c^2 + 2ac = 2s^2 + 2(56) = 2s^2 + 112$. Similarly, $(b+d)^2 = 2s^2 + 180$. 

 By Ptolemy's Theorem on $PCDA$, $as + cs = ds\sqrt{2}$, and therefore $a + c = d\sqrt{2}$. By Ptolemy's on $PBAD$, $bs + ds = as\sqrt{2}$, and therefore $b + d = a\sqrt{2}$. By squaring both equations, we obtain
 \begin{alignat*}{8} 2d\textasciicircum{}2 &= (a+c)\textasciicircum{}2 &&= 2s\textasciicircum{}2 + 112, \\ 2a\textasciicircum{}2 &= (b+d)\textasciicircum{}2 &&= 2s\textasciicircum{}2 + 180. \end{alignat*}Thus, $a^2 = s^2 + 90$, and $d^2 = s^2 + 56$. Plugging these values into $a^2 + c^2 = b^2 + d^2 = 2s^2$, we obtain $c^2 = s^2 - 90$, and $b^2 = s^2 - 56$. Now, we can solve using $a$ and $c$ (though using $b$ and $d$ yields the same solution for $s$).
 \begin{align*} ac = (\sqrt{s^2 - 90})(\sqrt{s^2 + 90}) &= 56 \\ (s^2 + 90)(s^2 - 90) &= 56^2 \\ s^4 &= 90^2 + 56^2 = 106^2 \\ s^2 &= \boxed{106}. \end{align*}
~mathboy100

 



\textbf{Solution 2 (Areas and Pythagorean Theorem)}

By the \textbf{Inscribed Angle Theorem}, we conclude that $\triangle PAC$ and $\triangle PBD$ are right triangles.

 Let the brackets denote areas. We are given that
 \begin{alignat*}{8} 2[PAC] &= PA \cdot PC &&= 56, \\ 2[PBD] &= PB \cdot PD &&= 90. \end{alignat*}Let $O$ be the center of the circle, $X$ be the foot of the perpendicular from $P$ to $\overline{AC},$ and $Y$ be the foot of the perpendicular from $P$ to $\overline{BD},$ as shown below:
                    \begin{figure}[H]
                        \centering
                        \includegraphics[width=0.5\textwidth]{Solution/AIME_I_2023_5_Solution_Figure_1.png}
                    \end{figure}
                    Let $d$ be the diameter of $\odot O.$ It follows that 
 \begin{alignat*}{8} 2[PAC] &= d\cdot PX &&= 56, \\ 2[PBD] &= d\cdot PY &&= 90. \end{alignat*}Moreover, note that $OXPY$ is a rectangle. By the Pythagorean Theorem, we have 
 \[PX^2+PY^2=PO^2.\]We rewrite this equation in terms of $d:$ 
 \[\left(\frac{56}{d}\right)^2+\left(\frac{90}{d}\right)^2=\left(\frac d2\right)^2,\] from which $d^2=212.$ Therefore, we get 
 \[[ABCD] = \frac{d^2}{2} = \boxed{106}.\]~MRENTHUSIASM

 



\textbf{Solution 3 (Similar Triangles)}


                    \begin{figure}[H]
                        \centering
                        \includegraphics[width=0.5\textwidth]{Solution/AIME_I_2023_5_Solution_Figure_1.png}
                    \end{figure}
                    Let the center of the circle be $O$, and the radius of the circle be $r$. Since $ABCD$ is a rhombus with diagonals $2r$ and $2r$, its area is $\dfrac{1}{2}(2r)(2r) = 2r^2$. Since $AC$ and $BD$ are diameters of the circle, $\triangle APC$ and $\triangle BPD$ are right triangles. Let $X$ and $Y$ be the foot of the altitudes to $AC$ and $BD$, respectively. We have
 \[[\triangle APC] = \frac{1}{2}(PA)(PC) = \frac{1}{2}(PX)(AC),\]so $PX = \dfrac{(PA)(PC)}{AC} = \dfrac{28}{r}$. Similarly,
 \[[\triangle BPD] = \frac{1}{2}(PB)(PD) = \frac{1}{2}(PY)(PB),\]so $PY = \dfrac{(PB)(PD)}{BD} = \dfrac{45}{r}$. Since $\triangle APX \sim \triangle PCX,$
 \[\frac{AX}{PX} = \frac{PX}{CX}\]
 \[\frac{AO - XO}{PX} = \frac{PX}{OC + XO}.\]But $PXOY$ is a rectangle, so $PY = XO$, and our equation becomes
 \[\frac{r - PY}{PX} = \frac{PX}{r + PY}.\]Cross multiplying and rearranging gives us $r^2 = PX^2 + PY^2 = \left(\dfrac{28}{r}\right)^2 + \left(\dfrac{45}{r}\right)^2$, which rearranges to $r^4 = 2809$. Therefore $[ABCD] = 2r^2 = \boxed{106}$.

 



\textbf{Solution 4 (Heights and Half-Angle Formula)}

Drop a height from point $P$ to line $\overline{AC}$ and line $\overline{BC}$. Call these two points to be $X$ and $Y$, respectively. Notice that the intersection of the diagonals of $\square ABCD$ meets at a right angle at the center of the circumcircle, call this intersection point $O$.

 Since $OXPY$ is a rectangle, $OX$ is the distance from $P$ to line $\overline{BD}$. We know that $\tan{\angle{POX}} = \frac{PX}{XO} = \frac{28}{45}$ by triangle area and given information. Then, notice that the measure of $\angle{OCP}$ is half of $\angle{XOP}$.

 Using the half-angle formula for tangent,

 
 \begin{align*} \frac{(2 \cdot \tan{\angle{OCP}})}{(1-\tan^2{\angle{OCP}})} = \tan{\angle{POX}} = \frac{28}{45} \\ 14\tan^2{\angle{OCP}} + 45\tan{\angle{OCP}} - 14 = 0 \end{align*}


 Solving the equation above, we get that $\tan{\angle{OCP}} = -7/2$ or $2/7$. Since this value must be positive, we pick $\frac{2}{7}$. Then, $\frac{PA}{PC} = 2/7$ (since $\triangle CAP$ is a right triangle with line $\overline{AC}$ the diameter of the circumcircle) and $PA * PC = 56$. Solving we get $PA = 4$, $PC = 14$, giving us a diagonal of length $\sqrt{212}$ and area $\boxed{106}$.

 



\textbf{Solution 5 (Analytic Geometry)}

Denote by $x$ the half length of each side of the square.We put the square to the coordinate plane, with $A = \left( x, x \right)$, $B = \left( - x , x \right)$, $C = \left( - x , - x \right)$, $D = \left( x , - x \right)$.

 The radius of the circumcircle of $ABCD$ is $\sqrt{2} x$.Denote by $\theta$ the argument of point $P$ on the circle.Thus, the coordinates of $P$ are $P = \left( \sqrt{2} x \cos \theta , \sqrt{2} x \sin \theta \right)$.

 Thus, the equations $PA \cdot PC = 56$ and $PB \cdot PD = 90$ can be written as
 \begin{align*} \sqrt{\left( \sqrt{2} x \cos \theta - x \right)^2 + \left( \sqrt{2} x \sin \theta - x \right)^2} \cdot \sqrt{\left( \sqrt{2} x \cos \theta + x \right)^2 + \left( \sqrt{2} x \sin \theta + x \right)^2} & = 56 \\ \sqrt{\left( \sqrt{2} x \cos \theta + x \right)^2 + \left( \sqrt{2} x \sin \theta - x \right)^2} \cdot \sqrt{\left( \sqrt{2} x \cos \theta - x \right)^2 + \left( \sqrt{2} x \sin \theta + x \right)^2} & = 90 \end{align*}


 These equations can be reformulated as
 \begin{align*} x^4 \left( 4 - 2 \sqrt{2} \left( \cos \theta + \sin \theta \right) \right) \left( 4 + 2 \sqrt{2} \left( \cos \theta + \sin \theta \right) \right) & = 56^2  \\ x^4 \left( 4 + 2 \sqrt{2} \left( \cos \theta - \sin \theta \right) \right) \left( 4 - 2 \sqrt{2} \left( \cos \theta - \sin \theta \right) \right) & = 90^2 \end{align*}


 These equations can be reformulated as
 \begin{align*} 2 x^4 \left( 1 - 2 \cos \theta  \sin \theta \right) & = 28^2 \hspace{1cm} (1) \\ 2 x^4 \left( 1 + 2 \cos \theta  \sin \theta \right) & = 45^2 \hspace{1cm} (2) \end{align*}


 Taking $\frac{(1)}{(2)}$, by solving the equation, we get
 \[ 2 \cos \theta \sin \theta = \frac{45^2 - 28^2}{45^2 + 28^2} . \hspace{1cm} (3) \]

 Plugging (3) into (1), we get
 \begin{align*} {\textrm Area} \ ABCD & = \left( 2 x \right)^2 \\ & = 4 \sqrt{\frac{28^2}{2 \left( 1 - 2 \cos \theta \sin \theta \right)}} \\ & = 2 \sqrt{45^2 + 28^2} \\ & = 2 \cdot 53 \\ & = \boxed{\textbf{(106) }} . \end{align*}


 



\textbf{Solution 6 (Law of Cosines)}

WLOG, let $P$ be on minor arc $\overarc {AB}$. Let $r$ and $O$ be the radius and center of the circumcircle respectively, and let $\theta = \angle AOP$.  	

 By the \href{https://artofproblemsolving.com/wiki/index.php?title=Pythagorean\_Theorem}{Pythagorean Theorem}, the area of the square is $2r^2$. We can use the \href{https://artofproblemsolving.com/wiki/index.php?title=Law\_of\_Cosines}{Law of Cosines} on isosceles triangles $\triangle AOP, \, \triangle COP, \, \triangle BOP, \, \triangle DOP$ to get 

 
 \begin{align*} 	 PA^2 &= 2r^2(1 - \cos \theta), \\	 PC^2 &= 2r^2(1 - \cos (180  - \theta)) = 2r^2(1 + \cos \theta), \\	 PB^2 &= 2r^2(1 - \cos (90 - \theta)) = 2r^2(1 - \sin \theta), \\	 PD^2 &= 2r^2(1 - \cos (90 + \theta)) = 2r^2(1 + \sin \theta).	 \end{align*}


 Taking the products of the first two and last two equations, respectively, 
 \[56^2 = (PA \cdot PC)^2 = 4r^4(1 - \cos \theta)(1 + \cos \theta) = 4r^4(1 - \cos^2 \theta) = 4r^4 \sin^2 \theta,\] and 
 \[90^2 = (PB \cdot PD)^2 = 4r^4(1 - \sin \theta)(1 + \sin \theta) = 4r^4(1 - \sin^2 \theta) = 4r^4 \cos^2 \theta.\]	Adding these equations, 
 \[56^2 + 90^2 = 4r^4,\] so  	
 \[2r^2 = \sqrt{56^2+90^2} = 2\sqrt{28^2+45^2} = 2\sqrt{2809} = 2 \cdot 53 = \boxed{106}.\]	~OrangeQuail9

 



\textbf{Solution 7 (Subtended Chords)}

First draw a diagram.
                    \begin{figure}[H]
                        \centering
                        \includegraphics[width=0.5\textwidth]{Solution/AIME_I_2023_5_Solution_Figure_3.png}
                    \end{figure}
                    Let's say that the radius is $r$. Then the area of the $ABCD$ is $(\sqrt2r)^2 = 2r^2$Using the formula for the length of a chord subtended by an angle, we get
 \[PA = 2r\sin\left(\dfrac{\theta}2\right)\]
 \[PC = 2r\sin\left(\dfrac{180-\theta}2\right) = 2r\sin\left(90 - \dfrac{\theta}2\right) = 2r\cos\left(\dfrac{\theta}2\right)\]Multiplying and simplifying these 2 equations gives
 \[PA \cdot PC = 4r^2 \sin \left(\dfrac{\theta}2 \right) \cos \left(\dfrac{\theta}2 \right) = 2r^2 \sin\left(\theta \right) = 56\]Similarly $PB = 2r\sin\left(\dfrac{90 +\theta}2\right)$ and $PD =2r\sin\left(\dfrac{90 -\theta}2\right)$. Again, multiplying gives
 \[PB \cdot PD = 4r^2 \sin\left(\dfrac{90 +\theta}2\right) \sin\left(\dfrac{90 -\theta}2\right) = 4r^2 \sin\left(90 -\dfrac{90 -\theta}2\right) \sin\left(\dfrac{90 -\theta}2\right)\]
 \[=4r^2 \sin\left(\dfrac{90 -\theta}2\right) \cos\left(\dfrac{90 -\theta}2\right) = 2r^2 \sin\left(90 - \theta \right) = 2r^2 \cos\left(\theta \right) = 90\]Dividing $2r^2 \sin \left(\theta \right)$ by $2r^2 \cos \left( \theta \right)$ gives $\tan \left(\theta \right) = \dfrac{28}{45}$, so $\theta = \tan^{-1} \left(\dfrac{28}{45} \right)$.Pluging this back into one of the equations, gives
 \[2r^2 = \dfrac{90}{\cos\left(\tan^{-1}\left(\dfrac{28}{45}\right)\right)}\]If we imagine a $28$-$45$-$53$ right triangle, we see that if $28$ is opposite and $45$ is adjacent, $\cos\left(\theta\right) = \dfrac{\text{adj}}{\text{hyp}} = \dfrac{45}{53}$. Now we see that
 \[2r^2 = \dfrac{90}{\frac{45}{53}} = \boxed{106}.\]~Voldemort101

 



\textbf{Solution 8 (Coordinates and Algebraic Manipulation)}


                    \begin{figure}[H]
                        \centering
                        \includegraphics[width=0.5\textwidth]{Solution/AIME_I_2023_5_Solution_Figure_4.png}
                    \end{figure}
                    Let $P=(a,b)$ on the upper quarter of the circle, and let $k$ be the side length of the square. Hence, we want to find $k^2$. Let the center of the circle be $(0,0)$.The two equations would thus become:
 \[\left(\left(a+\dfrac{k}2\right)^2+\left(b-\dfrac{k}2\right)^2\right)\left(\left(a-\dfrac{k}2\right)^2+\left(b+\dfrac{k}2\right)^2\right)=56^2\]
 \[\left(\left(a-\dfrac{k}2\right)^2+\left(b-\dfrac{k}2\right)^2\right)\left(\left(a+\dfrac{k}2\right)^2+\left(b+\dfrac{k}2\right)^2\right)=90^2\]Now, let $m=\left(a+\dfrac{k}2\right)^2$, $n=\left(a-\dfrac{k}2\right)^2$, $o=\left(b+\dfrac{k}2\right)^2$, and $p=\left(b-\dfrac{k}2\right)^2$. Our equations now change to $(m+p)(n+o)=56^2=mn+op+mo+pn$ and $(n+p)(m+o)=90^2=mn+op+no+pm$. Subtracting the first from the second, we have $pm+no-mo-pn=p(m-n)-o(m-n)=(m-n)(p-o)=34\cdot146$. Substituting back in and expanding, we have $2ak\cdot-2bk=34\cdot146$, so $abk^2=-17\cdot73$. We now have one of our terms we need ($k^2$). Therefore, we only need to find $ab$ to find $k^2$.We now write the equation of the circle, which point $P$ satisfies: 
 \[a^2+b^2=\left(\dfrac{k\sqrt{2}}{2}\right)^2=\dfrac{k^2}2\]We can expand the second equation, yielding 
 \[\left(a^2+b^2+\dfrac{k^2}2+(ak+bk)\right)\left(a^2+b^2+\dfrac{k^2}2-(ak+bk)\right)=(k^2+k(a+b))(k^2-k(a+b))=8100.\]Now, with difference of squares, we get $k^4-k^2\cdot(a+b)^2=k^2(k^2-(a+b)^2)=8100$. We can add $2abk^2=-17\cdot73\cdot2=-2482$ to this equation, which we can factor into $k^2(k^2-(a+b)^2+2ab)=k^2(k^2-(a^2+b^2))=8100-2482$. We realize that $a^2+b^2$ is the same as the equation of the circle, so we plug its equation in: $k^2\left(k^2-\dfrac{k^2}2\right)=5618$. We can combine like terms to get $k^2\cdot\dfrac{k^2}2=5618$, so $(k^2)^2=11236$.Since the answer is an integer, we know $11236$ is a perfect square. Since it is even, it is divisible by $4$, so we can factor $11236=2^2\cdot2809$. With some testing with approximations and last-digit methods, we can find that $53^2=2809$. Therefore, taking the square root, we find that $k^2$, the area of square $ABCD$, is $2\cdot53=\boxed{106}$.

 



\textbf{Solution 9 (Law of Sines)}

WLOG, let $P$ be on minor arc $AD.$ Draw in $AP$, $BP$, $CP$, $DP$ and let $\angle ABP = x.$  We can see, by the inscribed angle theorem, that $\angle APB = \angle ACB = 45$, and $\angle CPD = \angle CAD = 45.$ Then, $\angle PAB = 135-x$, $\angle PCD = \angle PAD = (135-x)-90 = 45-x$, and $\angle PDC = 90+x.$ Letting $(PA, PB, PC, PD, AB) = (a,b,c,d,s)$, we can use the law of sines on triangles $PAB$ and $PCD$ to get 
 \[s\sqrt{2} = \frac{a}{\sin(x)} = \frac{b}{\sin(135-x)} = \frac{c}{\sin(90+x)} = \frac{d}{\sin(45-x)}.\] Making all the angles in the above equation acute gives 
 \[s\sqrt{2} = \frac{a}{\sin(x)} = \frac{b}{\sin(45+x)} = \frac{c}{\sin(90-x)} = \frac{d}{\sin(45-x)}.\] 

 Note that we are looking for $s^{2}.$ We are given that $ac = 56$ and $bd = 90.$ This means that$s^{2}\sin(x)\sin(90-x) = 28$ and $s^{2}\sin(45+x)\sin(45-x) = 45.$ However, 
 \[\sin(x)\sin(90-x) = \sin(x)\cos(x) = \frac{\sin(2x)}{2}\] and 
 \[\sin(45+x)\sin(45-x) = \frac{(\cos(x) + \sin(x))(\cos(x) - \sin(x))}{2} = \frac{\cos^{2}(x) - \sin^{2}(x)}{2} = \frac{\cos(2x)}{2}.\]  Therefore, $s^{2}\sin(2x) = 56$ and $s^{2}\cos(2x) = 90.$ Therefore, by the Pythagorean Identity, 
 \[s^{2} = \sqrt{(s^{2}\sin(2x))^{2} + (s^{2}\cos(2x))^{2}} = \sqrt{56^{2} + 90^{2}} = \boxed{106}.\]

 



\textbf{Solution 10 (Areas and Trigonometry)}

Similar to Solution 6, let $P$ be on minor arc $\overarc {AB}$, $r$ and $O$ be the radius and center of the circumcircle respectively, and $\theta = \angle AOP$. Since $\triangle APC$ is a right triangle, $PA \cdot PC$ equals the hypotenuse, $2r$, times its altitude, which can be represented as $r \sin \theta$. Therefore, $2r^2 \sin \theta = 56$. Applying similar logic to $\triangle BPD$, we get $2r^2 \sin (90^\circ - \theta) = 2r^2 \cos \theta = 90$.

 Dividing the two equations, we have
 \begin{align*} \frac{\sin \theta}{\cos \theta} &= \frac{56}{90} \\ 56 \cos \theta &= 90 \sin \theta \\ (56 \cos \theta)^2 &= (90 \sin \theta)^2. \end{align*}
Adding $(56 \sin \theta)^2$ to both sides allows us to get rid of $\cos \theta$:
 \begin{align*} (56 \cos \theta)^2 + (56 \sin \theta)^2 &= (90 \sin \theta)^2 + (56 \sin \theta)^2 \\ 56^2 &= (90^2 + 56^2)(\sin \theta)^2 \\ \frac{56^2}{90^2 + 56^2} &= (\sin \theta)^2 \\ \frac{28}{53} &= \sin \theta. \end{align*}
Therefore, we have $2r^2\left(\frac{28}{53}\right) = 56$, and since the area of the square can be represented as $2r^2$, the answer is $56 \cdot \frac{53}{28} = \boxed{106}$.

 


 



\textbf{Solution 11 (Angle Chasing and Trigonometric Identities)}

First, we define a few points. Let $O$ be the center of the circle, let $E$ be the intersection of diameter $AC$ and chord $PD$, and let $F$ be the intersection of diameter $BD$ and chord $PC$. We know that $A$, $B$, $C$, and $D$ are four corners of a square. Therefore, the arcs $AD$, $DC$, and $CB$ are all $90$ degrees. By inscribed angles, angle $APD$, angle $DPC$, and angle $CPB$ are $45$ degrees each. Let the measure of angle $PAC$ be $a$. Similarly, let the measure of angle $PBD$ be $b$. 

 
Angle chasing will lead us to the fact that $a + b = 135$, or rather, $b = 135-a$. Let the diameter of the circle be $d$. Given by the problem, $d^2\sin a \cos a = 56$. Also, $d^2\sin b \cos b = 90$. Using the trigonometric identity $\sin 2x = 2\sin x \cos x$, we can rewrite these as $d^2\sin 2a = 112$ and $d^2\sin 2b = 180$. Since we determined that $b = \frac{3\pi}{4}-a$, this can be substituted into the second equation. Then, we divide the two equations to get $\frac{\sin (\frac{3\pi}{2}-2a)}{\sin 2a} = \frac{45}{28}$. By using the trigonometric difference-of-angle identity, this simplifies to $\frac{-\cos 2a}{\sin 2a} = \frac{45}{28}$. By the definition of the tangent function, $\tan 2a = -\frac{28}{45}$

 
Considering this hypothetical right triangle with legs of $28$ and $45$, the hypotenuse is $\sqrt{45^2+28^2} = 53$. Since $\sin 2a$ must be positive (since $a$ is acute), $\sin 2a = \frac{28}{53}$. Substituting this into the first of the equations, $\frac{28}{53}d^2 = 112$. From this, $d^2 = 212$. The area of square $ABCD$ is half of the square of its diagonal, which is $d$. Thus, the answer is $\frac{d^2}{2} = \boxed{106}$.

 



\textbf{Solution 12 (Six Consecutive Equations)}

This solution is like solution 1, but with more algebra than actual geometry.

 Firstly, like in solution 1, label the respective lengths $PA$, $PB$, $PC$, $PD$, $AB = BC = CD = DA$ as $a$, $b$, $c$, $d$, $s$.

 While only five equations is required, we use six to apply the shortcut of solving multistep equations, elimination.

 Our six equations are as desired.

 \begin{align}(5.1)\quad & ac = 56 && \text{Given} \\(5.2)\quad & bd = 90 && \text{Given} \\(5.3)\quad & b^2 + d^2 = 2s^2 && \text{Pythagorean on } \triangle BPD \\(5.4)\quad & a^2 + c^2 = 2s^2 && \text{Pythagorean on } \triangle APC \\(5.5)\quad & bs + as\sqrt{2} = ds && \text{Ptolemy on } APBD \\(5.6)\quad & bs + ds = cs\sqrt{2} && \text{Ptolemy on } CBPD\end{align}

 Note that rearranging $(5.6)$ gives $(5.6R)~bs - cs\sqrt{2} = -ds$ and thus $(5.5) + (5.6R)$ results in $2bs + as\sqrt{2} - cs\sqrt{2} = 0$, or $2b + \sqrt{2}(a-c) = 0$, and thus $b = \frac{\sqrt{2}(c-a)}{2}$. It then becomes obvious that $b$'s companion $d$ is $d = \frac{\sqrt{2}(c+a)}{2}$. Thus $(5.2)$ can be arranged to become $\frac{1}{2} \cdot (c^2 - a^2) = 90$, and thus $c^2 -  a^2 = 180$. We have that $ab = 56$ and $c^2 - a^2 = 180$, in which we can then find $a$ and $c$ to be $4$ and $14$ (discard negative solutions as negative lengths do not exist in Plane Euclidean Geometry) respectively. Thus, we have for our final equation $a^2 + c^2 = 2s^2$, and we just need the result $s^2$ to solve the problem. Then our answer is simply $\frac{196+16}{2} = \frac{212}{2} = \boxed{106}$. $\square$

 



\textbf{Solution 13 (Heavy coord bash)}

Let the center of the square/circle be the origin, and assume WLOG that $P$ is on the arc connecting $A$ to $B$. Let $A=\left(\dfrac{-r\sqrt{2}}{2},\dfrac{r\sqrt{2}}{2}\right)$, and the other coordinates go in clockwise order. Finally, let $P=(x,y)$, with $x^2+y^2=r^2$.

 From the first equation, $(PA)(PC)=56$, we have\begin{align*}\sqrt{\left(\left(x+\dfrac{r\sqrt2}{2}\right)^2+\left(y-\dfrac{r\sqrt2}{2}\right)^2\right)\left(\left(x-\dfrac{r\sqrt2}{2}\right)^2+\left(y+\dfrac{r\sqrt2}{2}\right)^2\right)}&=56\\\sqrt{(x^2+y^2+xr\sqrt2-yr\sqrt2+r^2)(x^2+y^2-xr\sqrt2+yr\sqrt2+r^2)}&=56\\\sqrt{(2r^2+(xr\sqrt2-yr\sqrt2))(2r^2-(xr\sqrt2-yr\sqrt2))}&=56\\\sqrt{4r^4-(xr\sqrt2-yr\sqrt2)^2}&=56\\\sqrt{4r^4-(2x^2r^2+2y^2r^2-2xyr^2\sqrt2)}&=56\\\sqrt{2r^4+2xyr^2\sqrt2}&=56\\2r^4+2xyr^2\sqrt2&=56^2.\end{align*}
Using $(PB)(PD)=90$, we get a very similar equation:
 \[2r^4-2xyr^2\sqrt2=90^2.\]Adding the equations, we have
 \[4r^4=90^2+56^2.\]Dividing both sides by $4$, we have $r^4=45^2+28^2$, Since $28,45,$ and $53$ form a Pythagorean Triple, we have $r^2=53$. Since the side length of the square is $r\sqrt2$, the area is thus $2r^2=\boxed{106}.$

 


	\item (\textit{AIME I 2023}) Alice knows that $3$ red cards and $3$ black cards will be revealed to her one at a time in random order. Before each card is revealed, Alice must guess its color. If Alice plays optimally, the expected number of cards she will guess correctly is $\frac{m}{n},$ where $m$ and $n$ are relatively prime positive integers. Find $m+n.$


 
 \textbf{Answer: }051



\textbf{Solution 1 (Casework)}

We break the problem into stages, one for each card revealed, then further into cases based on the number of remaining unrevealed cards of each color. Since \href{https://artofproblemsolving.com/wiki/index.php?title=Expected\_value}{expected value} is linear, the expected value of the total number of correct card color guesses across all stages is the sum of the expected values of the number of correct card color guesses at each stage; that is, we add the probabilities of correctly guessing the color at each stage to get the final answer (See \href{https://brilliant.org/wiki/linearity-of-expectation/}{https://brilliant.org/wiki/linearity-of-expectation/})

 At any stage, if there are $a$ unrevealed cards of one color and $b$ of the other color, and $a \geq b$, then the optimal strategy is to guess the color with $a$ unrevealed cards, which succeeds with probability $\frac{a}{a+b}.$

 \textbf{Stage 1:} 

 There are always $3$ unrevealed cards of each color, so the probability of guessing correctly is $\frac{1}{2}$.

 \textbf{Stage 2:}

 There is always a $3$-$2$ split ($3$ unrevealed cards of one color and $2$ of the other color), so the probability of guessing correctly is $\frac{3}{5}$.

 \textbf{Stage 3:}

 There are now $2$ cases:

 
                \begin{itemize}
                    \item  The guess from Stage 2 was correct, so there is now a $2$-$2$ split of cards and a $\frac{1}{2}$ probability of guessing the color of the third card correctly.\item  The guess from Stage 2 was incorrect, so the split is $3$-$1$ and the probability of guessing correctly is $\frac{3}{4}$. 
                \end{itemize}


                Thus, the overall probability of guessing correctly is $\frac{3}{5} \cdot \frac{1}{2} + \frac{2}{5} \cdot \frac{3}{4} = \frac{3}{5}$.

 \textbf{Stage 4:}

 This stage has $2$ cases as well:

 
                \begin{itemize}
                    \item  The guesses from both Stage 2 and Stage 3 were incorrect. This occurs with probability $\frac{2}{5} \cdot \frac{1}{4} = \frac{1}{10}$ and results in a $3$-$0$ split and a certain correct guess at this stage.\item  Otherwise, there must be a $2$-$1$ split and a $\frac{2}{3}$ probability of guessing correctly.
                \end{itemize}


                The probability of guessing the fourth card correctly is therefore $\frac{1}{10} \cdot 1 + \frac{9}{10} \cdot \frac{2}{3} = \frac{7}{10}$.

 \textbf{Stage 5:}

 Yet again, there are $2$ cases:

 
                \begin{itemize}
                    \item  In Stage 4, there was a $2$-$1$ split and the guess was correct. This occurs with probability $\frac{9}{10} \cdot \frac{2}{3} = \frac{3}{5}$ and results in a $1$-$1$ split with a $\frac{1}{2}$ chance of a correct guess here.\item  Otherwise, there must be a $2$-$0$ split, making a correct guess certain.
                \end{itemize}


                In total, the fifth card can be guessed correctly with probability $\frac{3}{5} \cdot \frac{1}{2} + \frac{2}{5} \cdot 1 = \frac{7}{10}$.

 \textbf{Stage 6:}

 At this point, only $1$ card remains, so the probability of guessing its color correctly is $1$.

 In conclusion, the expected value of the number of cards guessed correctly is 
 \[\frac{1}{2} + \frac{3}{5} + \frac{3}{5} + \frac{7}{10} + \frac{7}{10} + 1 = \frac{5+6+6+7+7+10}{10} = \frac{41}{10},\] so the answer is $41 + 10 = \boxed{051}.$

 



\textbf{Solution 2 (Casework)}

At any point in the game, Alice should guess whichever color has come up less frequently thus far (although if both colors have come up equally often, she may guess whichever she likes); using this strategy, her probability of guessing correctly is at least $\frac{1}{2}$ on any given card, as desired.

 There are ${6 \choose 3} = 20$ possible orderings of cards, all equally likely (since any of the $6! = 720$ permutations of the cards is equally likely, and each ordering covers $3!^2 = 6^2 = 36$ permutations).

 Each of the $10$ orderings that start with red cards corresponds with one that starts with a black card; the problem is symmetrical with respect to red and black cards, so we can, without loss of generality, consider only the orderings that start with red cards.

 We then generate a tally table showing whether Alice's guesses are correct for each ordering; for a given card, she guesses correctly if fewer than half the previously shown cards were the same color, guesses incorrectly if more than half were the same color, and guesses correctly with probability $\frac{1}{2}$ if exactly half were the same color. 

 In this table, $\mid$ denotes a correct guess, $\--$ denotes an incorrect guess, and $/$ denotes a guess with $\frac{1}{2}$ probability of being correct.

 
                    \begin{figure}[H]
                        \centering
                        \includegraphics[width=0.5\textwidth]{Solution/AIME_I_2023_6_Solution_Figure_1.png}
                    \end{figure}
                    Now we sum the tallies across orderings, obtaining $41$, and finally divide by the number of orderings ($10$) to obtain the expected number of correct guesses, $\frac{41}{10}$, which yields an answer of $41 + 10 = \boxed{051}.$

 



\textbf{Solution 3 (Dynamic Programming)}

Denote by $N \left( a, b \right)$ the optimal expected number of cards that Alice guesses correctly, where the number of red and black cards are $a$ and $b$, respectively.

 Thus, for $a, b \geq 1$, we have
 \begin{align*} N \left( a, b \right) & = \max \left\{ \frac{a}{a+b} \left( 1 + N \left( a - 1 , b \right) \right) + \frac{b}{a+b} N \left( a , b - 1 \right) , \right. \\ & \hspace{1cm} \left. \frac{a}{a+b} N \left( a - 1 , b \right) + \frac{b}{a+b} \left( 1 + N \left( a , b - 1 \right) \right) \right\} . \end{align*}


 For $a = 0$, Alice always guesses black. So $N \left( 0 , b \right) = b$.

 For $b = 0$, Alice always guesses red. So $N \left( a , 0 \right) = a$.

 To solve this dynamic program, we can also exploit its symmetry that $N \left( a , b \right) = N \left( b , a \right)$.

 By solving this dynamic program, we get$N \left( 1, 1 \right) = \frac{3}{2}$, $N \left( 1, 2 \right) = \frac{7}{3}$, $N \left( 1 , 3 \right) = \frac{13}{4}$, $N \left( 2 , 2 \right) = \frac{17}{6}$, $N \left( 2 , 3 \right) = \frac{18}{5}$, $N \left( 3, 3 \right) = \frac{41}{10}$.

 Therefore, the answer is $41 + 10 = \boxed{\textbf{(051) }}$.

 



\textbf{Solution 4 (Simplification of Solution 3)}

Denote by $N_{i,j}$ the optimal expected number of cards that Alice guesses correctly, where the number of cards are $i$ and $j \ge i.$

 If $i = 0$ then Alice guesses correctly all cards, so $N_{0,j} = j.$

 If $j = i$ then Alice guesses next card with probability  $\frac {1}{2} \implies N_{i,i} = \frac {1}{2} + N_{i-1,i}.$

 If $j = i+1$ then Alice guesses next card with probability $\frac {i+1}{2i+1} \implies N_{i,i+1} = \frac {i+1}{2i+1} (1+ N_{i,i}) + \frac{i}{2i+1}  N_{i-1,i+1}.$

 If $j = i+2$ then Alice guesses next card with probability  $\frac {i+2}{2i+2} \implies N_{i,i+2} = \frac {i+2}{2i+2} (1+ N_{i,i+1}) + \frac{i}{2i+2}  N_{i-1,i+2}.$

 One can find consistently: $N_{1,1} = \frac {1}{2} + N_{0,1} = \frac {3}{2},$
 \[N_{1,2} = \frac {2}{3} (1 + N_{1,1}) + \frac {1}{3} N_{0,2} = \frac {7}{3}.\]
 \[N_{2,2} = \frac {1}{2} + N_{1,2} = \frac {17}{6}.\]
 \[N_{1,3} = \frac {3}{4} (1 + N_{1,2}) + \frac {1}{4} N_{0,3} = \frac {13}{4}.\]
 \[N_{2,3} = \frac {3}{5} (1 + N_{2,2}) + \frac {2}{5} N_{1,3} = \frac {18}{5}.\]
 \[N_{3,3} = \frac {1}{2} + N_{2,3} = \frac {41}{10}.\]Therefore, the answer is $41 + 10 = \boxed{\textbf{(051) }}$.

 \textbf{vladimir.shelomovskii@gmail.com, vvsss}

 



\textbf{Solution 5 (Pseudo-recursion)}

Denote $E_n$ the expected number of cards Alice guesses correctly given $n$ red cards and $n$ black cards. We want to find $E_3$.

 Alice has a $\frac{1}{2}$ chance of guessing the first card. WLOG assume the first card color is red. For the next card, Alice has a $\frac{3}{5}$ chance of guessing the card if she chooses black; if they guess right, there's one less red and black card, so the expected number of cards Alice guesses from here is $E_2$. If Alice does not guess correctly (which occurs with probability $\frac{2}{5}$), this means that there's 3 black cards and 1 red card left, so Alice should guess black next with a $\frac{3}{4}$ chance of being right. If Alice is wrong (with probability $\frac{1}{4}$), there are only 3 black cards left, so Alice can guess these with certainty; if Alice is right, there are 2 blacks and 1 red left, so Alice should again guess black. If Alice is right (with probability $\frac{2}{3}$), there is now 1 black and red card each, so the expected number of cards guessed is $E_1$; if she is wrong (with probability $\frac{1}{3}$), there are 2 black cards left, so Alice can guess these with certainty.

 Summing this up into a formula:
 \[E_3 = \frac{1}{2} + \frac{3}{5} \left(1 + E_2 \right) + \frac{2}{5} \left( \frac{1}{4}(3) + \frac{3}{4}\left(1 + \frac{2}{3}\left(1 + E_1 \right) + \frac{1}{3}(2)\right) \right)\]

 We can apply similar logic to compute $E_2$ and get
 \[E_2 = \frac{1}{2} + \frac{2}{3}(1 + E_1) + \frac{1}{3}(2)\]

 To compute $E_1$, we know that Alice can guess the last card with certainty, and there's a $\frac{1}{2}$ chance they get the first card as well, so $E_1 = \frac{3}{2}$.

 Thus, $E_2 = \frac{17}{6}$, and after long computation, we get $E_3 = \frac{41}{10}$. The requested answer is $41 + 10 = \boxed{51}$.

 



\textbf{Solution 6 (non-recursive)}

There's a beautiful observation that for any particular game path, the expected number of correct guesses is purely dependent on how often we are in a neutral state (i.e. there are the same number of red and black cards remaining in the deck).  For example, given path $BRRBBR$ we are in a neutral state once in the beginning when no cards have been flipped, after the second card has been flipped, after the 4th card has been flipped, and finally a 4th time at the very end of the game.  Call the number of times a given path is in the neutral state $n$.  Then the expected number of correct guesses equals:

 
 \[E = 3 + \frac{n-1}{2}\]

 Since this problem is small, we could just could just count them (noting that the problem is symmetric between reds and blacks):

 
 \[ \begin{array}{|c|c|c|} \hline \textbf{Path} & \textbf{Neutral States} & \textbf{E}\\ \hline \text{RRRBBB} & 2 & 3.5\\ \text{RRBRBB} & 2 & 3.5 \\ \text{RRBRBB} & 3 & 4 \\ \text{RRBBBR} & 3 & 4 \\ \text{RBRRBB} & 3 & 4 \\ \text{RBRBRB} & 4 & 4.5 \\ \text{RBRBBR} & 4 & 4.5 \\ \text{RBBRRB} & 4 & 4.5 \\ \text{RBBRBR} & 4 & 4.5 \\ \text{RBBBRR} & 3 & 4 \\ \hline \end{array} \]

 Using linearity of expectation, the expected number of correct guesses is:

 
 \[E=\frac{3.5+3.5+4+4+4+4.5+4.5+4.5+4.5+4}{10}=\frac{41}{10}, \text{ thus } 41 + 10 = \boxed{51}\]

 Note that if we call the total number of neutral states across all ${6 \choose 3}=20$ paths $N$, this would be equivalent to:

 
 \[E=3 + \frac{\frac{N}{20}-1}{2}\]

 We can compute $N$ directly.  Let the index $m$ be a point where we've seen $m$ red cards and $m$ black cards.  There are $\binom{2m}{m}$ paths that lead up to that point. After which there are $6-2m$ cards remaining, $3-m$ of which are red and $3-m$ of which are black.  Summing up over all possible values of $m$ gives us:

 
 \[ \sum_{m=0}^{3} \binom{2m}{m} \cdot \binom{6-2m}{3-m} =  \] \[ \binom{0}{0} \cdot \binom{6}{3} + \binom{2}{1} \cdot \binom{4}{2} + \binom{4}{2} \cdot \binom{2}{1} + \binom{6}{3} \cdot \binom{0}{0} \] \[ = 1 \cdot 20 + 2 \cdot 6 + 6 \cdot 2 + 20 \cdot 1  \] \[ = 20 + 12 + 12 + 20 = 64 \]

 and as before:

 
 \[E=3 + \frac{\frac{N}{20}-1}{2}=3 + \frac{\frac{64}{20}-1}{2}=3+\frac{11}{10}=\frac{41}{10}\]

 


	\item (\textit{AIME I 2023}) Call a positive integer $n$ extra-distinct if the remainders when $n$ is divided by $2, 3, 4, 5,$ and $6$ are distinct. Find the number of extra-distinct positive integers less than $1000$.


 
 \textbf{Answer: }049



\textbf{Solution 1}

$n$ can either be $0$ or $1$ (mod $2$).

 Case 1: $n \equiv 0 \pmod{2}$

 Then, $n \equiv 2 \pmod{4}$, which implies $n \equiv 1 \pmod{3}$ and $n \equiv 4 \pmod{6}$, and therefore $n \equiv 3 \pmod{5}$. Using \href{https://artofproblemsolving.com/wiki/index.php?title=Chinese\_Remainder\_Theorem}{CRT}, we obtain $n \equiv 58 \pmod{60}$, which gives $16$ values for $n$.

 Case 2: $n \equiv 1 \pmod{2}$

 $n$ is then $3 \pmod{4}$. If $n \equiv 0 \pmod{3}$, $n \equiv 3 \pmod{6}$, a contradiction. Thus, $n \equiv 2 \pmod{3}$, which implies $n \equiv 5 \pmod{6}$. $n$ can either be $0 \pmod{5}$, which implies that $n \equiv 35 \pmod{60}$ by CRT, giving $17$ cases; or $4 \pmod{5}$, which implies that $n \equiv 59 \pmod{60}$ by CRT, giving $16$ cases.

 The total number of extra-distinct numbers is thus $16 + 16 + 17 = \boxed{049}$.

 



\textbf{Solution 2 (Simpler)}

Because the LCM of all of the numbers we are dividing by is $60$, we know that all of the remainders are $0$ again at $60$, meaning that we have a cycle that repeats itself every $60$ numbers.

 After listing all of the remainders up to $60$, we find that $35$, $58$, and $59$ are extra-distinct. So, we have $3$ numbers every $60$ which are extra-distinct. $60\cdot16 = 960$ and $3\cdot16 = 48$, so we have $48$ extra-distinct numbers in the first $960$ numbers. Because of our pattern, we know that the numbers from $961$ thru $1000$ will have the same remainders as $1$ thru $40$, so we have $1$ other extra-distinct number ($35$). 

 $48 + 1 =  \boxed{049}$.

 



\textbf{Solution 3}

$\textbf{Case 0:  } {\textrm Rem} \ \left( n, 6 \right) = 0$.

 We have ${\textrm Rem} \ \left( n, 2 \right) = 0$. This violates the condition that $n$ is extra-distinct.Therefore, this case has no solution.

 $\textbf{Case 1:  } {\textrm Rem} \ \left( n, 6 \right) = 1$.

 We have ${\textrm Rem} \ \left( n, 2 \right) = 1$. This violates the condition that $n$ is extra-distinct.Therefore, this case has no solution.

 $\textbf{Case 2:  } {\textrm Rem} \ \left( n, 6 \right) = 2$.

 We have ${\textrm Rem} \ \left( n, 3 \right) = 2$. This violates the condition that $n$ is extra-distinct.Therefore, this case has no solution.

 $\textbf{Case 3:  } {\textrm Rem} \ \left( n, 6 \right) = 3$.

 The condition ${\textrm Rem} \ \left( n, 6 \right) = 3$ implies ${\textrm Rem} \ \left( n, 2 \right) = 1$, ${\textrm Rem} \ \left( n, 3 \right) = 0$.

 Because $n$ is extra-distinct, ${\textrm Rem} \ \left( n, l \right)$ for $l \in \left\{ 2, 3, 4 \right\}$ is a permutation of $\left\{ 0, 1 ,2 \right\}$.Thus, ${\textrm Rem} \ \left( n, 4 \right) = 2$.

 However, ${\textrm Rem} \ \left( n, 4 \right) = 2$ conflicts ${\textrm Rem} \ \left( n, 2 \right) = 1$.Therefore, this case has no solutions.

 $\textbf{Case 4:  } {\textrm Rem} \ \left( n, 6 \right) = 4$.

 The condition ${\textrm Rem} \ \left( n, 6 \right) = 4$ implies ${\textrm Rem} \ \left( n, 2 \right) = 0$ and ${\textrm Rem} \ \left( n, 3 \right) = 1$.

 Because $n$ is extra-distinct, ${\textrm Rem} \ \left( n, l \right)$ for $l \in \left\{ 2, 3, 4 , 5 \right\}$ is a permutation of $\left\{ 0, 1 ,2 , 3 \right\}$.

 Because ${\textrm Rem} \ \left( n, 2 \right) = 0$, we must have ${\textrm Rem} \ \left( n, 4 \right) = 2$. Hence, ${\textrm Rem} \ \left( n, 5 \right) = 3$.

 Hence, $n \equiv -2 \pmod{{\textrm lcm} \left( 4, 5 , 6 \right)}$.Hence, $n \equiv - 2 \pmod{60}$.

 We have $1000 = 60 \cdot 16 + 40$.Therefore, the number of extra-distinct $n$ in this case is 16.

 $\textbf{Case 5:  } {\textrm Rem} \ \left( n, 6 \right) = 5$.

 The condition ${\textrm Rem} \ \left( n, 6 \right) = 5$ implies ${\textrm Rem} \ \left( n, 2 \right) = 1$ and ${\textrm Rem} \ \left( n, 3 \right) = 2$.

 Because $n$ is extra-distinct, ${\textrm Rem} \ \left( n, 4 \right)$ and ${\textrm Rem} \ \left( n, 5 \right)$ are two distinct numbers in $\left\{ 0, 3, 4 \right\}$.Because ${\textrm Rem} \ \left( n, 4 \right) \leq 3$ and $n$ is odd, we have ${\textrm Rem} \ \left( n, 4 \right) = 3$.Hence, ${\textrm Rem} \ \left( n, 5 \right) = 0$ or 4.

 $\textbf{Case 5.1:  } {\textrm Rem} \ \left( n, 6 \right) = 5$, ${\textrm Rem} \ \left( n, 4 \right) = 3$, ${\textrm Rem} \ \left( n, 5 \right) = 0$.

 We have $n \equiv 35 \pmod{60}$.

 We have $1000 = 60 \cdot 16 + 40$.Therefore, the number of extra-distinct $n$ in this subcase is 17.

 $\textbf{Case 5.2:  } {\textrm Rem} \ \left( n, 6 \right) = 5$, ${\textrm Rem} \ \left( n, 4 \right) = 3$, ${\textrm Rem} \ \left( n, 5 \right) = 4$.

 $n \equiv - 1 \pmod{60}$.

 We have $1000 = 60 \cdot 16 + 40$.Therefore, the number of extra-distinct $n$ in this subcase is 16.

 Putting all of the cases together, the total number of extra-distinct $n$ is $16 + 17 + 16 = \boxed{049}$.

 



\textbf{Solution 4 (Small addition to solution 2)}

We need to find that $35$, $58$, and $59$ are all extra-distinct numbers smaller than $61.$

 Let $k \in \left\{2, 3, 4, 5, 6 \right\}.$ Denote the remainder in the division of $a$ by $b$ as ${\textrm Rem} \ \left( a, b \right).$

 ${\textrm Rem} \ \left( -1, k \right) = k - 1 \implies {\textrm Rem} \ \left( 59, k \right) = k - 1 = \left\{1, 2, 3, 4, 5 \right\}\implies 59$ is extra-distinct.${\textrm Rem} \ \left( -2, k \right) = k - 2 \implies {\textrm Rem} \ \left( 58, k \right) = k - 2 = \left\{0, 1, 2, 3, 4 \right\} \implies 58$ is extra-distinct.
 \[{\textrm Rem} \ \left( x + 12y, k \right) = {\textrm Rem} \ \left(x, k \right) + \left\{0, 0, 0, {\textrm Rem} \ \left(12y, k \right), 0 \right\}.\]We need to check all of the remainders up to $12 - 3 = 9$ and remainders 
 \[{\textrm Rem} \ \left( 59 - 12, k \right) = {\textrm Rem} \ \left( 59 - 36, k \right) = \left\{1, 2, 3, 3, 5 \right\}, {\textrm Rem} \ \left( 59 - 48, k \right) = \left\{1, 2, 3, 1, 5 \right\},\]${\textrm Rem} \ \left( 59 - 24, k \right) ={\textrm Rem} \ \left(35, k \right) = \left\{1, 2, 3, 0, 5 \right\} \implies 35$ is extra-distinct.$58 - 12 = 46 \implies {\textrm Rem} \ \left( 46, 5 \right) = 1 = {\textrm Rem} \ \left( 46, 3 \right),$ $58 - 24 = 34 \implies {\textrm Rem} \ \left( 34, 5 \right) = 4 = {\textrm Rem} \ \left( 34, 6 \right),$ $58 - 36 = 22 \implies {\textrm Rem} \ \left( 22, 5 \right) = 2 = {\textrm Rem} \ \left( 22, 4 \right),$ $58 - 48 = 10 \implies {\textrm Rem} \ \left( 10, 5 \right) = 0 = {\textrm Rem} \ \left( 10, 2 \right).$

 \textbf{vvsss}

 


	\item (\textit{AIME I 2023}) Rhombus $ABCD$ has $\angle BAD < 90^\circ.$ There is a point $P$ on the incircle of the rhombus such that the distances from $P$ to the lines $DA,AB,$ and $BC$ are $9,5,$ and $16,$ respectively. Find the perimeter of $ABCD.$


 
 \textbf{Answer: }125



\textbf{Solution 1}

This solution refers to the \textbf{Diagram} section.

 Let $O$ be the incenter of $ABCD$ for which $\odot O$ is tangent to $\overline{DA},\overline{AB},$ and $\overline{BC}$ at $X,Y,$ and $Z,$ respectively. Moreover, suppose that $R,S,$ and $T$ are the feet of the perpendiculars from $P$ to $\overleftrightarrow{DA},\overleftrightarrow{AB},$ and $\overleftrightarrow{BC},$ respectively, such that $\overline{RT}$ intersects $\odot O$ at $P$ and $Q.$

 We obtain the following diagram:
                    \begin{figure}[H]
                        \centering
                        \includegraphics[width=0.5\textwidth]{Solution/AIME_I_2023_8_Solution_Figure_1.png}
                    \end{figure}
                    Note that $\angle RXZ = \angle TZX = 90^\circ$ by the properties of tangents, so $RTZX$ is a rectangle. It follows that the diameter of $\odot O$ is $XZ = RT = 25.$

 Let $x=PQ$ and $y=RX=TZ.$ We apply the Power of a Point Theorem to $R$ and $T:$
 \begin{align*} y^2 &= 9(9+x), \\ y^2 &= 16(16-x). \end{align*}
We solve this system of equations to get $x=7$ and $y=12.$ Alternatively, we can find these results by the symmetry on rectangle $RTZX$ and semicircle $\widehat{XPZ}.$

 We extend $\overline{SP}$ beyond $P$ to intersect $\odot O$ and $\overleftrightarrow{CD}$ at $E$ and $F,$ respectively, where $E\neq P.$ So, we have $EF=SP=5$ and $PE=25-SP-EF=15.$ On the other hand, we have $PX=15$ by the Pythagorean Theorem on right $\triangle PRX.$ Together, we conclude that $E=X.$ Therefore, points $S,P,$ and $X$ must be collinear.

 Let $G$ be the foot of the perpendicular from $D$ to $\overline{AB}.$ Note that $\overline{DG}\parallel\overline{XP},$ as shown below:
                    \begin{figure}[H]
                        \centering
                        \includegraphics[width=0.5\textwidth]{Solution/AIME_I_2023_8_Solution_Figure_2.png}
                    \end{figure}
                    As $\angle PRX = \angle AGD = 90^\circ$ and $\angle PXR = \angle ADG,$ we conclude that $\triangle PRX \sim \triangle AGD$ by the AA Similarity. The ratio of similitude is 
 \[\frac{PX}{AD} = \frac{RX}{GD}.\] We get $\frac{15}{AD} = \frac{12}{25},$ from which $AD = \frac{125}{4}.$

 Finally, the perimeter of $ABCD$ is $4AD = \boxed{125}.$

 



\textbf{Solution 2}


                    \begin{figure}[H]
                        \centering
                        \includegraphics[width=0.5\textwidth]{Solution/AIME_I_2023_8_Solution_Figure_3.png}
                    \end{figure}
                    Let $G$ be the extension of line $CD$ such that $CG \perp FG,$ $E$ be the point where $AB$ is tangent to the circle, $H$ be where the perpendicular from $F$ to line $GC$ meets the circle after point $P$, and $I$ be the point where $PH$ meets $AD$ (point $I$ is added to make it clear that $H$ is not on $AD$). 

 Connect $H$ to the point where the circle is tangent to $CD$, call this point $K.$ Since $EK$ is a diameter, $PH$ is a chord, and $EK \parallel PH$ (they are parallel because $EFGK$ is a rectangle), the perpendicular bisector of $PH$ will also perpendicularly bisect $EK.$ So trapezoid $PHKE$ is symmetrical about this perpendicular bisector and is therefore isosceles. Then $PE = HK.$ So by Hypotenuse-Leg congruence $\triangle EFP \cong \triangle KGH.$ Then $HG = FP = 9$ and since the length of $FG$ is $25,$ $PH = 7.$  Now, consider the right triangle formed by the radius, height, and half of the top base of the isosceles trapezoid. Let $h$ be the height. We have $h^2 + (7/2)^2 = (25/2)^2 \implies h = FE = 12.$ Then by the Pythagorean Theorem on $\triangle EFP$, $EP = 15.$

 
We observe that $\angle FPE \cong \angle EAJ$ since right triangles $EFP$ and $EJA$ share common angle $E.$

 Therefore, $\sin(\angle FPE) = \frac{12}{15} = \frac{4}{5} = \sin(\angle EAJ)$.

 The height of the rhombus is $25$ since the diameter of the circle is $25$. If $s$ is the side length of the rhombus, $h$ is the height, and $\theta$ is an interior angle, then  $\sin(\theta) = \frac{h}{s}$ (you can see this by dropping an altitude from $B$ to $AD$).

 Given $\sin(\theta) = \frac{4}{5}$, we have $\frac{25}{s} = \frac{4}{5}$.

 Solving, we find $s = \frac{125}{4}$.

 So the perimeter of the rhombus is $4s = 4 \cdot \frac{125}{4} = \boxed{125}$.

 



\textbf{Solution 3}

This solution refers to the \textbf{Diagram} section.

 Define points $O,R,S,$ and $T$ as Solution 1 does. Moreover, let $H$ be the foot of the perpendicular from $P$ to $\overleftrightarrow{CD},$ $M$ be the foot of the perpendicular from $O$ to $\overleftrightarrow{HS},$ and$N$ be the foot of the perpendicular from $O$ to $\overleftrightarrow{RT}.$

 We obtain the following diagram:
                    \begin{figure}[H]
                        \centering
                        \includegraphics[width=0.5\textwidth]{Solution/AIME_I_2023_8_Solution_Figure_4.png}
                    \end{figure}
                    Note that the diameter of $\odot O$ is $HS=RT=25,$ so $OP=\frac{25}{2}.$ It follows that:

 <ol style="margin-left: 1.5em;"> \item In right $\triangle OMP,$ we have $MP=\frac{HS}{2}-PS=\frac{15}{2}$ by symmetry, from which $OM=10$ by the Pythagorean Theorem. 

 \item In right $\triangle ONP,$ we have $NP=\frac{RT}{2}-RP=\frac{7}{2}$ by symmetry, from which $ON=12$ by the Pythagorean Theorem.

 </ol>Since $\overline{MO}\parallel\overline{AB}$ and $\overline{ON}\parallel\overline{DA},$ we conclude that $\angle A = \angle MON.$ We apply the Sine of a Sum Formula:
 \begin{align*} \sin\angle A &= \sin\angle MON \\ &= \sin(\angle MOP + \angle PON) \\ &= \sin\angle MOP \cos\angle PON + \cos\angle MOP \sin\angle PON \\ &= \frac{3}{5}\cdot\frac{24}{25} + \frac{4}{5}\cdot\frac{7}{25} \\ &= \frac{4}{5}. \end{align*}
Note that 
 \[\sin\angle A = \frac{HS}{DA},\] from which $\frac{4}{5} = \frac{25}{DA}.$ We solve this equation to get $DA=\frac{125}{4}.$

 Finally, the perimeter of $ABCD$ is $4DA = \boxed{125}.$

 



\textbf{Solution 4}

Label the points of the rhombus to be $X$, $Y$, $Z$, and $W$ and the center of the incircle to be $O$ so that $9$, $5$, and $16$ are the distances from point $P$ to side $ZW$, side $WX$, and $XY$ respectively. Through this, we know that the distance from the two pairs of opposite lines of rhombus $XYZW$ is $25$ and circle $O$ has radius $\frac{25}{2}$. 

 Call the feet of the altitudes from $P$ to side $ZW$, side $WX$, and side $XY$ to be $A$, $B$, and $C$ respectively. Additionally, call the feet of the altitudes from $O$ to side $ZW$, side $WX$, and side $XY$ to be $D$, $E$, and $F$ respectively.

 Draw a line segment from $P$ to $\overline{OD}$ so that it is perpendicular to $\overline{OD}$. Notice that this segment length is equal to $AD$ and is $\sqrt{\left(\frac{25}{2}\right)^2-\left(\frac{7}{2}\right)^2}=12$ by Pythagorean Theorem.

 Similarly, perform the same operations with perpendicular from $P$ to $\overline{OE}$ to get $BE=10$.

 By equal tangents, $WD=WE$. Now, label the length of segment $WA=n$ and $WB=n+2$.

 Using Pythagorean Theorem again, we get

 
 \begin{align*} WA^2+PA^2&=WB^2+PB^2 \\ n^2+9^2&=(n+2)^2+5^2 \\ n&=13. \end{align*}


 Which also gives us $\tan{\angle{OWX}}=\frac{1}{2}$ and $OW=\frac{25\sqrt{5}}{2}$.

 Since the diagonals of the rhombus intersect at $O$ and are angle bisectors and are also perpendicular to each other, we can get that

 
 \begin{align*} \frac{OX}{OW}&=\tan{\angle{OWX}} \\ OX&=\frac{25\sqrt{5}}{4} \\ WX^2&=OW^2+OX^2 \\ WX&=\frac{125}{4} \\ 4WX&=\boxed{125}. \end{align*}


 



\textbf{Solution 5}

Denote by $O$ the center of $ABCD$.We drop an altitude from $O$ to $AB$ that meets $AB$ at point $H$.We drop altitudes from $P$ to $AB$ and $AD$ that meet $AB$ and $AD$ at $E$ and $F$, respectively.We denote $\theta = \angle BAC$.We denote the side length of $ABCD$ as $d$.

 Because the distances from $P$ to $BC$ and $AD$ are $16$ and $9$, respectively, and $BC \parallel AD$, the distance between each pair of two parallel sides of $ABCD$ is $16 + 9 = 25$.Thus, $OH = \frac{25}{2}$ and $d \sin \theta = 25$.

 We have
 \begin{align*} \angle BOH & = 90^\circ - \angle HBO \\ & = 90^\circ - \angle HBD \\ & = 90^\circ - \frac{180^\circ - \angle C}{2} \\ & = 90^\circ - \frac{180^\circ - \theta}{2} \\ & = \frac{\theta}{2} . \end{align*}


 Thus, $BH = OH \tan \angle BOH = \frac{25}{2} \tan \frac{\theta}{2}$.

 In $FAEP$, we have $\overrightarrow{FA} + \overrightarrow{AE} + \overrightarrow{EP} + \overrightarrow{PF} = 0$.Thus,
 \[ AF + AE e^{i \left( \pi - \theta \right)} + EP e^{i \left( \frac{3 \pi}{2} - \theta \right)} - PF i . \]

 Taking the imaginary part of this equation and plugging $EP = 5$ and $PF = 9$ into this equation, we get
 \[ AE = \frac{9 + 5 \cos \theta}{\sin \theta} . \]

 We have
 \begin{align*} OP^2 & = \left( OH - EP \right)^2 + \left( AH - AE \right)^2 \\ & = \left( \frac{25}{2} - 5 \right)^2 + \left( d - \frac{25}{2} \tan \frac{\theta}{2} - \frac{9 + 5 \cos \theta}{\sin \theta} \right) \\ & = \left( \frac{15}{2} \right)^2 + \left( \frac{25}{\sin \theta} - \frac{25}{2} \tan \frac{\theta}{2} - \frac{9 + 5 \cos \theta}{\sin \theta} \right) . \hspace{1cm} (\bigstar) \end{align*}


 Because $P$ is on the incircle of $ABCD$, $OP = \frac{25}{2}$. Plugging this into $(\bigstar)$, we get the following equation
 \[ 20 \sin \theta - 15 \cos \theta = 7 . \]

 By solving this equation, we get $\sin \theta = \frac{4}{5}$ and $\cos \theta = \frac{3}{5}$.Therefore, $d = \frac{25}{\sin \theta} = \frac{125}{4}$.

 Therefore, the perimeter of $ABCD$ is $4d = \boxed{125}$.

 



\textbf{Solution 6}

The center of the incircle is $O.$ Denote the points in which the incircle meets $\overline{AB},$ $\overline{BC},$ $\overline{CD},$ and $\overline{DA}$ as $W,$ $X,$ $Y,$ and $Z,$ respectively. Next, also denote the base of the perpendicular from $P$ to $\overline{AB},$ $\overline{AD},$ $\overline{OW},$ and $\overline{OZ}$ as $M,$ $N,$ $S,$ and $T,$ respectively.

 We can easily see that the radius of the circle is $\frac{25}{2}.$ Using this and Pythagorus on right $\triangle OSP$ and $\triangle OTP,$ we find that $MW = PS = 10$ and $NZ = PT = 12.$

 Since $AW = AZ$ by properties of circle tangents, we can deduce by the above information that $AM = AN+2.$ Doing Pythagorus on right $\triangle AMP$ and $\triangle ANP$ we find that $a^2 = b^2 + 56$ (because $a^2+25=b^2+81.$) From solving the $2$ just derived equations, we find that $AM=15$ and $AN=13.$

 Next, we use Pythagorus on right $\triangle AOB$ (we can see it's right because of properties of rhombuses.) We get 
 \[AB^2 = AO^2 + BO^2.\] We know $AB = AW + WB = 25 + WB.$ By Pythagorus on $\triangle AWO$ and $\triangle BWO,$ we also know $AO^2 = 25^2+\left(\frac{25}{2}\right)^2$ and $BO^2=WB^2+\left(\frac{25}{2}\right)^2.$ Substituting these in, we have 
 \[25^2 + 50WB + WB^2 = 25^2+\left(\frac{25}{2}\right)^2+\left(\frac{25}{2}\right)^2+WB^2.\] Solving for $WB,$ we get $WB = \frac{25}{4}.$ Now we find that each side of the rhombus $=AB=25+\frac{25}{4}=\frac{125}{4}.$ The perimeter of the rhombus would be that times $4.$ Our final answer is 
 \[\frac{125}{4}\cdot4=\boxed{125}.\]

 



\textbf{Solution 7}


                    \begin{figure}[H]
                        \centering
                        \includegraphics[width=0.5\textwidth]{Solution/AIME_I_2023_8_Solution_Figure_5.png}
                    \end{figure}
                    Notation is shown on diagram, $RT \perp AD, FG \perp AB, E = AD \cap \omega, E' = FG \cap AD.$$RT = 9 + 16 = 25 = FG$ as hights of rhombus.
 \[RP = QT = 9, PQ = 16 - 9 = 7, GE' = PF = 5,\]
 \[PE' = 25 - 5 - 5 = 15, RE = \sqrt{RP \cdot RQ} = \sqrt{9 \cdot 16} = 12.\]
 \[PE =  \sqrt{RP^2 + RE^2} = 15 \implies E = E'.\]
 \[\sin \alpha = \frac {RE}{PE} =  \frac {GF}{AD} \implies AD = \frac {15 \cdot 25}{12} = \frac {125}{4}.\]The perimeter of $ABCD$ is $\frac{125}{4}\cdot4=\boxed{125}.$

 \textbf{vladimir.shelomovskii@gmail.com, vvsss}

 


	\item (\textit{AIME I 2023}) Find the number of cubic polynomials $p(x) = x^3 + ax^2 + bx + c,$ where $a, b,$ and $c$ are integers in $\{-20,-19,-18,\ldots,18,19,20\},$ such that there is a unique integer $m \not= 2$ with $p(m) = p(2).$


 
 \textbf{Answer: }738



\textbf{Solution 1}

Plugging $2$ and $m$ into the expression for $p(x)$ and equating them, we get $8+4a+2b+c = m^3+am^2+bm+c$. Rearranging, we have 
 \[(m^3-8) + (m^2 - 4)a + (m-2)b = 0.\] Note that the value of $c$ won't matter as it can be anything in the provided range, giving a total of $41$ possible choices for $c.$ So we only need to find the number of ordered pairs $(a, b)$ that work, and multiply it by $41.$We can start by first dividing both sides by $m-2.$ (Note that this is valid since $m\neq2:$ 
 \[m^2 + 2m + 4 + (m+2)a + b = 0.\] We can rearrange this so it is a quadratic in $m$: 
 \[m^2 + (a+2)m + (4 + 2a + b) = 0.\] Remember that $m$ has to be unique and not equal to $2.$ We can split this into two cases: case $1$ being that $m$ has exactly one solution, and it isn't equal to $2$; case $2$ being that $m$ has two solutions, one being $2,$ but the other is a unique solution not equal to $2.$

 
$\textbf{Case 1:}$

 There is exactly one solution for $m,$ and that solution is not $2.$ This means that the discriminant of the quadratic equation is $0,$ using that, we have $(a+2)^2 = 4(4 + 2a + b)$. Rearranging, we have 
 \[(a-2)^2 = 4(4 + b)\Longrightarrow a = 2\pm2\sqrt{4+b}.\] Using the fact that $4+b$ must be a perfect square, we can easily see that the values for $b$ can be $-4, -3, 0, 5,$ and $12.$ Also since it's a "$\pm$" there will usually be $2$ solutions for $a$ for each value of $b.$ The two exceptions for this would be if $b = -4$ and $b = 12.$ For $b=-4$ because it would be a $\pm0,$ which only gives one solution, instead of two. For $b=12$, because then $a = -6$, the solution for $m$ would equal $2,$, and we don't want this. (We can know this by replacing the solutions into the quadratic formula). 

 So we have $5$ solutions for $b,$ each of which gives $2$ values for $a,$ except for $2$ and $b = -4,$ both of which only give one. So in total, there are $5*2 - 2 = 8$ ordered pairs of $(a,b)$ in this case.

 
$\textbf{Case 2:}$

 $m$ has two solutions, but exactly one of them isn't equal to $2.$ This ensures that $1$ of the solutions is equal to $2.$

 Let $r$ be the other value of $m$ that isn't $2.$ By Vieta: 
 \begin{align*} r+2 &= -a-2\\ 2r &= 4+2a+b. \end{align*}
 From the first equation, we subtract both sides by $2$ and double both sides to get $2r = -2a - 8$ which also equals to $4+2a+b$ from the second equation. Equating both, we have $4a + b + 12 = 0.$ We can easily count that there would be $11$ ordered pairs $(a,b)$ that satisfy that.

 However, there's an outlier case in which $r$ happens to also equal $2,$ and we don't want that. We can reverse engineer and find out that $r=2$ when $(a,b) = (-6, 12),$ which we overcounted. So we subtract by one and we conclude that there are $10$ ordered pairs of $(a,b)$ that satisfy this case.

 
This all shows that there is a total of $8+10 = 18$ amount of ordered pairs $(a,b).$ Multiplying this by $41$ (the amount of values for $c$) we get $18\cdot41=\boxed{738}$ as our final answer.

 



\textbf{Solution 2 (factor the difference)}

$p(x)-p(2)$ is a cubic with at least two integral real roots, therefore it has three real roots, which are all integers.

 There are exactly two distinct roots, so either $p(x)=p(2)+(x-2)^2(x-m)$ or $p(x)=p(2)+(x-2)(x-m)^2$, with $m\neq 2$.

 In the first case $p(x)=x^3-(4+m)x^2+(4+4m)x-4m+p(2)$, with $|4+4m|\leq 20$ (which entails $|4+m|\leq 20$), so $m$ can be $-6,-5,-4,-3,-2,-1,0,1, (\textbf{not 2}!), 3,4$ and $-4m+p(2)$ can be any integer from $-20$ to $20$, giving $410$ polynomials (since the coefficients are given by linear functions of $m$ and thus are distinct). 

 In the second case $p(x)=x^3-(2+2m)x^2+(4m+m^2)x-2m^2+p(2)$, and $m$ can be $-6,-5,-4,-3,-2,-1,0,1$ and $-4m+p(2)$ can be any integer from $-20$ to $20$, giving $328$ polynomials.

 The total is $\boxed{738}$.

 



\textbf{Solution 3 (derivative)}

For the equation $p(x) = p(2)$ to have exactly two distinct real roots, the line $y = p(2)$ must be tangent to either the local maximum or the local minimum of the cubic. This occurs when the derivative $p'(x) = 0$ at one of the roots.

 Case 1: The repeated root is at $x=2$.

 If $x=2$ is a repeated root, then $p'(2) = 0$. Given $p(x) = x^3 + ax^2 + bx + c$, the derivative is $p'(x) = 3x^2 + 2ax + b$. $p'(2) = 3(2)^2 + 2a(2) + b = 12 + 4a + b = 0 \implies b = -4a - 12.$ The constraint $|b| \le 20$ implies: $-20 \le -4a - 12 \le 20 \implies -8 \le a \le 2.$ This gives $11$ possible integers for $a$. However, we must exclude $a = -6$ because it makes $p''(2) = 6(2) + 2(-6) = 0$, creating an inflection point rather than an extremum (which would result in only one distinct root). Thus, there are $11 - 1 = 10$ valid pairs of $(a, b)$. Each pair combined with the $41$ possible values of $c$ gives $10 \times 41 = 410$ polynomials.

 Case 2: The repeated root is at $x=m$ ($m \neq 2$).

 In this case, the polynomial takes the form $p(x) - p(2) = (x-2)(x-m)^2$. Expanding this: $p(x) - p(2) = x^3 - (2m+2)x^2 + (m^2+4m)x - 2m^2.$ Matching the coefficients with $p(x) = x^3 + ax^2 + bx + c$, we have:

 $a = -2m - 2$

 $b = m^2 + 4m$

 For $a$ and $b$ to be integers, $m$ must be an integer. Testing values of $m$ such that $|a| \le 20$ and $|b| \le 20$: The condition $|a| \le 20$ gives $-11 \le m \le 9$. Within this range, checking $|m^2 + 4m| \le 20$ yields $m \in {-6, -5, -4, -3, -2, -1, 0, 1}$. Note that $m=2$ is excluded. This gives $8$ valid pairs of $(a, b)$. Combined with $41$ values of $c$, we have $8 \times 41 = 328$ polynomials.

 The total number of such polynomials is $410 + 328 = \boxed{738}.$

 


	\item (\textit{AIME I 2023}) There exists a unique positive integer $a$ for which the sum 
 \[U=\sum_{n=1}^{2023}\left\lfloor\dfrac{n^{2}-na}{5}\right\rfloor\] is an integer strictly between $-1000$ and $1000$. For that unique $a$, find $a+U$.


 (Note that $\lfloor x\rfloor$ denotes the greatest integer that is less than or equal to $x$.)


 
 \textbf{Answer: }944



\textbf{Solution (Super Simple, Easy, & Intuitive)}

Since 
 \[x = \left\lfloor x \right\rfloor + \{x\}\],\begin{align*}U &= \sum_{n=1}^{2023}\left\lfloor\frac{n^2-na}{5}\right\rfloor \\&= \frac15\left(\sum_{n=1}^{2023}n^2-a\sum_{n=1}^{2023}n\right) -\sum_{n=1}^{2023}\left\{\frac{n^2-na}{5}\right\} \\&= \frac{2023\cdot2024}{10}(1349-a) -\sum_{n=1}^{2023}\left\{\frac{n^2-na}{5}\right\}.\end{align*}


 
Since each fractional part is in $[0,1)$, the second sum is in $[0,2023)$. Thus, for $U\in(-1000,1000)$, we must have
 \[a=1349,\]which also proves uniqueness.

 Now $1349\equiv4\pmod5$, so

 


	\item (\textit{AIME I 2023}) Find the number of subsets of $\{1,2,3,\ldots,10\}$ that contain exactly one pair of consecutive integers. Examples of such subsets are $\{\mathbf{1},\mathbf{2},5\}$ and $\{1,3,\mathbf{6},\mathbf{7},10\}.$


 
 \textbf{Answer: }235



\textbf{Solution 1 (Minimal Casework)}

Define $f(x)$ to be the number of subsets of $\{1, 2, 3, 4, \ldots x\}$ that have $0$ consecutive element pairs, and $f'(x)$ to be the number of subsets that have $1$ consecutive pair.

 Using casework on where the consecutive element pair is, there is a unique consecutive element pair that satisfies the conditions. It is easy to see that 
 \[f'(10) = 2f(7) + 2f(6) + 2f(1)f(5) + 2f(2)f(4) + f(3)^2.\]

 We see that $f(1) = 2$, $f(2) = 3$, and $f(n) = f(n-1) + f(n-2)$. This is because if the element $n$ is included in our subset, then there are $f(n-2)$ possibilities for the rest of the elements (because $n-1$ cannot be used), and otherwise there are $f(n-1)$ possibilities. Thus, by induction, $f(n)$ is the $n+1$th Fibonacci number.

 This means that $f'(10) = 2(34) + 2(21) + 2(2)(13) + 2(3)(8) + 5^2 = \boxed{235}$.

 



\textbf{Solution 2}

We can solve this problem using casework, with one case for each possible pair of consecutive numbers.

 $\textbf{Case 1: (1,2)}$

 If we have (1,2) as our pair, we are left with the numbers from 3-10 as elements that can be added to our subset. So, we must compute how many ways we can pick these numbers so that the set has no consecutive numbers other than (1,2). Our first option is to pick no more numbers, giving us $8 \choose {0}$. We can also pick one number, giving us $7 \choose {1}$ because 3 cannot be picked. Another choice is to pick two numbers and in order to make sure they are not consecutive we must fix one number in between them, giving us $6 \choose {2}$. This pattern continues for each amount of numbers, yielding $5 \choose {3}$ for 3 numbers and $4 \choose {4}$ for four numbers. Adding these up, we have $8 \choose {0}$ + $7 \choose {1}$ + $6 \choose {2}$ + $5 \choose {3}$ + $4 \choose {4}$ = $\textbf{34}$.

 


 $\textbf{Case 2: (2,3)}$

 If we have (2,3) as our pair, everything works the same as with (1,2), because 1 is still unusable as it is consecutive with 2. The only difference is we now have only 4-10 to work with. Using the same pattern as before, we have $7 \choose {0}$ + $6 \choose {1}$ + $5 \choose {2}$ + $4 \choose {3}$ = $\textbf{21}$.

 


 $\textbf{Case 3: (3,4)}$

 This case remains pretty much the same except we now have an option of whether or not to include 1. If we want to represent this like we have with our other choices, we would say $2 \choose {0}$ for choosing no numbers and $1 \choose {1}$ for choosing 1, leaving us with $2 \choose {0}$ + $1 \choose {1}$ = 2 choices (either including the number 1 in our subset or not including it). As far as the numbers from 5-10, our pattern from previous cases still holds. We have $6 \choose {0}$ + $5 \choose {1}$ + $4 \choose {2}$ + $3 \choose {3}$ = 13. With 2 choices on one side and 13 choices on the other side, we have $2\cdot13$ = $\textbf{26}$ combinations in all.

 


 $\textbf{Case 4: (4,5)}$

 Following the patterns we have already created in our previous cases, for the numbers 1-3 we have $3 \choose {0}$ + $2 \choose {1}$ = 3 choices (1, 2, or neither) and for the numbers 6-10 we have $5 \choose {0}$ + $4 \choose {1}$ + $3 \choose {2}$ = 8 choices. With 3 choices on one side and 8 choices on the other side, we have $3\cdot8$ = $\textbf{24}$ combinations in all.

 


 $\textbf{Case 5: (5,6)}$

 Again following the patterns we have already created in our previous cases, for the numbers 1-4 we have $4 \choose {0}$ + $3 \choose {1}$ + $2 \choose {2}$ = 5 choices and for the numbers 5-10 we have the same $4 \choose {0}$ + $3 \choose {1}$ + $2 \choose {2}$ = 5 choices. $5\cdot5$ = $\textbf{25}$ combinations in all.

 


 $\textbf{Rest of the cases}$

 By symmetry, the case with (6,7) will act the same as case 4 with (4,5). This goes the same for (7,8) and case 3, (8,9) and case 2, and (9,10) and case 1.

 


 Now, we simply add up all of the possibilities for each case to get our final answer. 34 + 21 + 26 + 24 + 25 + 24 + 26 + 21 + 34 = $\boxed{\textbf{(235)}}$

 
-Algebraik

 



\textbf{Solution 3 (Double Recursive Equations)}

Denote by $N_1 \left( m \right)$ the number of subsets of a set $S$ that consists of $m$ consecutive integers, such that each subset contains exactly one pair of consecutive integers.

 Denote by $N_0 \left( m \right)$ the number of subsets of a set $S$ that consists of $m$ consecutive integers, such that each subset does not contain any consecutive integers.

 Denote by $a$ the smallest number in set $S$.

 First, we compute $N_1 \left( m \right)$.

 Consider $m \geq 3$.We do casework analysis.

 Case 1: A subset does not contain $a$.

 The number of subsets that has exactly one pair of consecutive integers is $N_1 \left( m - 1 \right)$.

 Case 2: A subset contains $a$ but does not contain $a + 1$.

 The number of subsets that has exactly one pair of consecutive integers is $N_1 \left( m - 2 \right)$.

 Case 3: A subset contains $a$ and $a + 1$.

 To have exactly one pair of consecutive integers, this subset cannot have $a + 2$, and cannot have consecutive integers in $\left\{ a+3, a+4, \cdots , a + m - 1 \right\}$.

 Thus, the number of subsets that has exactly one pair of consecutive integers is $N_0 \left( m - 3 \right)$.

 Therefore, for $m \geq 3$,
 \[N_1 \left( m \right) = N_1 \left( m - 1 \right) + N_1 \left( m - 2 \right) + N_0 \left( m - 3 \right) .\]

 For $m = 1$, we have $N_1 \left( 1 \right) = 0$.For $m = 2$, we have $N_1 \left( 2 \right) = 1$.

 
Second, we compute $N_0 \left( m \right)$.

 Consider $m \geq 2$.We do casework analysis.

 Case 1: A subset does not contain $a$.

 The number of subsets that has no consecutive integers is $N_0 \left( m - 1 \right)$.

 Case 2: A subset contains $a$.

 To avoid having consecutive integers, the subset cannot have $a + 1$.

 Thus, the number of subsets that has no consecutive integers is $N_0 \left( m - 2 \right)$.

 Therefore, for $m \geq 2$,
 \[N_0 \left( m \right) = N_0 \left( m - 1 \right) + N_0 \left( m - 2 \right)  .\]

 For $m = 0$, we have $N_0 \left( 0 \right) = 1$.For $m = 1$, we have $N_0 \left( 1 \right) = 2$.

 By solving the recursive equations above, we get $N_1 \left( 10 \right) = \boxed{\textbf{(235) }}$.

 



\textbf{Solution 4 (Recursive)}

Let $a_n$ be the number of cases from a subset ranging from $1$ - $n$.  We approach the subset from the start and the back to build up a equation. We figure out cases if $1$ and $10$ is in the subset.

 
If $1$ and $10$ are both not in the subset, We get $a_n$.

 
If $1$ is in the subset, and $10$ is not, We get multiple cases from this.If $2$ is in the subset, $3$ would not be in the subset. Now we have {$4$,$5$,$6$,$7$,$8$,$9$} and we'll have to determine how many subsets which there is no consecutive integers. After we do some testing, We get $a_n$ = $F_n+2$. From this case we get $F_8$.But if $2$ is not in the subset, We can just get $a_7$.The same goes for when $10$ is in the subset and $1$ is not, it is also $F_8$ + $a_7$.

 
But when $1$ and $10$ are both in the subset, We'll have to find more cases.If $2$ and $9$ are both not in the subset, We simply get $a_6$.When $2$ is in the subset and $9$ isn't, We know that $3$ is not in the subset since $1$ and $2$ forms a pair, The sum of this case will be $F_7$.The same also goes for when $9$ is in the subset and $2$ isnt.There is no cases for if $2$ and $9$ are both in the subset, Since that forms $2$ pairs.

 
Summarizing we get:$a_{10}$ = $a_8$ + $2$$a_7$ + $a_6$ + $2$$F_8$ + $2$$F_7$.

 
$2$$F_8$ + $2$$F_7$ = $2$$F_9$.

 
We get a equation of $a_{10}$ = $a_8$ + $2$$a_7$ + $a_6$ + $2$$F_9$

 We compute the basic cases.$a_1$ = $0$$a_2$ = $1$$a_3$ = $2$$a_4$ = $5$

 

 \[\begin{array}{rcl} a & F_a & Cases\\ \hline 1 & 1 & 0 \\ 2 & 1 & 1 \\ 3 & 2 & 2 \\ 4 & 3 & 5 \\ 5 & 5 & 10 \\ 6 & 8 & 20 \\ 7 & 13 & 38 \\ 8 & 21 & 71 \\ 9 & 34 & 130 \\ 10 & 55 & 235 \\ \end{array}\]

 
We get our answer,  $\boxed{235}$.

 



\textbf{Solution 5 (Similar to Solution 3)}

Let $a_n$ be the number of subsets of the set $\{1,2,3,\ldots,n\}$ such that there exists exactly 1 pair of consecutive elements. Let $b_n$ be the number of subsets of the set $\{1, 2, 3\ldots, n\}$ such that there doesn't exist any pair of consecutive elements. First, lets see how we can construct $a_n.$ For each subset $S$ counted in $a_n,$ either:1. $\{n-1, n\}\subseteq S,$2. $n\not\in S$, or3. $n-1 \not\in S$ and $n\in S.$The first case counts $b_{n-3}$ subsets (as $n-1$ cannot be included and the rest cannot have any consecutive elements), The second counts $a_{n-1},$ and the third counts $a_{n-2}.$ Thus, 
 \[a_n = a_{n-1} + a_{n-2} + b_{n-3}.\]Next, Lets try to construct $b_n.$ For each subset $T$ counted in $b_n,$ either:1.$n \not\in T,$ or 2.$n \in T.$The first case counts $b_{n-1}$ subsets and the second counts $b_{n-2}.$ Thus, 
 \[b_n = b_{n-1} + b_{n-2}.\]Since $b_1 = 2$ and $b_2 = 3,$ we have that $b_n = F_{n+1},$ so $a_n = a_{n-1} + a_{n-2} + F_{n-2}.$ (The $F_i$ is the $i$th Fibonacci number). From here, we can construct a table of the values of $a_n$ until $n = 10.$ By listing out possibilities, we can solve for our first 3 values. 

 
 \[\begin{array}{r|l} n & a_n \\ \hline 1 & 0 \\ 2 & 1\\ 3 & 2\\ 4 & 2 + 1 + F_2 = 5\\5 & 5 + 2 + F_3 = 10\\6 & 10 + 5 + F_4 = 20 \\ 7 & 20 + 10 + F_5 = 38\\ 8 & 38 + 20 + F_6 = 71 \\ 9& 71 + 38 + F_7 = 130\\ 10 & 130 + 71 + F_8 = 235 \end{array}\]Our answer is $a_{10} = \boxed{235}.$~AtharvNaphade

 



\textbf{Solution 6 (Stars and Bars)}

Note: This is a very common stars and bars application. 

 Casework on number of terms, let the number of terms be $n$. We can come up with a generalized formula for the number  of subsets with n terms.

 Let $d_1, d_2, ..., d_{n-1}$ be the differences between the n terms. \textit{For example, in the set {2, 3, 6}, $d_1 = 1; d_2 = 3; d_3 = 5$}

 Let the range of the set be k for now, $d_1 + ... + d_{n-1} = k$. We select one pair of terms to be consecutive by selecting one of the (n-1) terms to be 1. WLOG, let $d_1 = 1$. $1 + d_2 + ... + d_{n-1} = k$.

 To ensure the other $d_2, d_3, ..., d_{n-1}$ are greater than 1 such that no two other terms are consecutive or the same, let $D_2 = d_2 + 1; D_3 = d_3 + 1; ...$. $(n-2) + 1 + D_2 + ... + D_{n-1} = k$ where $D_2, D_3, ..., D_{n-1}$ are positive integers.

 Finally, we add in $D_0$, the distance between 0 and the first term of the set, and $D_n$, the distance between the last term and 11. This way, the "distance" from 0 to 11 is "bridged" by $D_0$, k, and $D_N$.
 \[D_0 + k + D_n = D_0 + ((n-1) + D_2 + ... + D_{n-1}) + D_N = 11\] 
 \[D_0 + D_2 + D_3 + ... + D_n = 12 - n\]There are n \textbf{positive terms}, by \href{https://artofproblemsolving.com/wiki/index.php?title=Balls\_and\_Urns}{Balls and Urns}, there are ${(12-n)-1 \choose (n)-1} = {11-n \choose n-1}$ ways of doing this. However, recall that there were $(n-1)$ ways, and we had used a WLOG to choose which two digits are consecutive. The final formula for the number of valid n-element subsets is hence $(n-1){11-n \choose n-1}$ for $n > 2$.

 Case 1: Two terms$n = 2$, so $(2-1){11-2 \choose 2-1} = 1{9 \choose 1} = 9$

 Case 2: Three terms$n = 3$, so $(3-1){11-3 \choose 3-1} = 2{8 \choose 2} = 56$

 Case 3: Four terms$n = 4$, so $(4-1){11-4 \choose 4-1} = 3{7 \choose 3} = 105$

 Case 4: Five terms$n = 5$, so $(5-1){11-5 \choose 5-1} = 4{6 \choose 4} = 60$

 Case 5: Six terms$n = 6$, so $(6-1){11-6 \choose 6-1} = 5{5 \choose 5} = 5$

 We can check by the Pigeonhole principle that there cannot be more than six terms, so the answer is $9+56+105+60+5=\boxed{235}$.

 



\textbf{Solution 7 (Fibonacci)}

Note that there are $F_{n+2}$ subsets of a set of $n$ consecutive integers that contains no two consecutive integers. (This can be proven by induction.)

 Now, notice that if we take $i$ and $i+1$ as the consecutive integers in our subset, we need to make a subset of the remaining integers such that it doesn't contain any two consecutive integers. Clearly, $i-1$ and $i+2$ cannot be chosen, and since $i-2$ and $i+3$ are sufficiently far apart, it is obvious we do not need to be concerned that an element of the set $\{1, 2, ... , i-2 \}$ is consecutive with any element of the set $\{ i+3, i+4, ... , 10 \}$ 

 Thus, we can count the number of ways to choose a subset from the first set without any two elements being consecutive and multiply this by the number of ways to choose a subset from the second set without any two elements being consecutive. From above, and noting that the first set has $i-2$ consecutive integer elements and the second set has $8-i$ consecutive integer elements, we know that this is $F_i F_{10-i}.$ 

 Summing this over for all $1 \leq i \leq 9$ yields 
 \[\sum_{i=1}^9 F_i F_{10-i} = F_1F_9 + F_2F_8 + F_3F_7 + F_4F_6 + F_5 F_5 + F_6 F_4 + F_7 F_3 + F_8 F_2 + F_9 F_1 = 2(F_1F_9 + F_2F_8 + F_3F_7+F_4F_6) + F_5^2 = 2(1 \cdot 34 + 1 \cdot 21 + 2 \cdot 13 + 3 \cdot 8) + 5^2 = 2 \cdot 105+25 = \boxed{235}.\]

 



\textbf{Solution 8 (Polyominoes)}

The problem is the same as laying out a line of  polynomoes to cover spots $0,1,...10$: 1 triomino ($RGG$), $n$ dominoes ($RG$), and $8-2n$ monominoes ($R$).  The $G$ spots cover the members of the subset.The total number spots is 11, because one $R$ spot always covers the 0, and the other spots cover 1 through 10.

 There are 5 ways to choose polyomino sets, and many ways to order each set:

 $R + RG + RGG =$ Polyominoes $\rightarrow$ Orderings 

 $0 + 4 + 1 = 5 \rightarrow 5! / 0!4!1! = ~~~5$

 $2 + 3 + 1 = 6 \rightarrow 6! / 2!3!1! = ~60$

 $4 + 2 + 1 = 7 \rightarrow 7! / 4!2!1! = 105$

 $6 + 1 + 1 = 8 \rightarrow 8! / 6!1!1! = ~56$

 $8 + 0 + 1 = 9 \rightarrow 9! / 8!0!1! = ~~~9$

 The sum is $\boxed{235}$.

 



\textbf{Solution 9 (Dynamic Programming)}

Let $dp(i)$ be the number of subsets of a set $i$ consecutive integers such that the maximum value in the set is $i$ and there exists exactly one pair of consecutive integers. Define $dp2(i)$ similarly, but without any pair of consecutive integers. The base cases are  $dp2(1)=dp2(2)=1$, $dp(1)=0$, and $dp(2)=1$.

 The transitions are:

 
 \[dp2(i)=\sum_{j=1}^{i-2}(dp2(j))+1\]

 
 \[dp(i)=dp2(i-1)+\sum_{j=1}^{i-2}dp(j)\]

 Note that $dp2$ is the Fibonacci numbers. 

 
 \[\begin{array}{rcl} i & dp & dp2\\ \hline 1 & 0 & 1 \\ 2 & 1 & 1 \\ 3 & 1 & 2 \\ 4 & 3 & 3 \\ 5 & 5 & 5 \\ 6 & 10 & 8 \\ 7 & 18 & 13 \\ 8 & 33 & 21 \\ 9 & 59 & 34 \\ 10 & 105 & 55 \\ \end{array}\]

 Summing over $dp$ yields $1+1+3+5+10+18+33+59+105=\boxed{235}$

 


	\item (\textit{AIME I 2023}) Let $\triangle ABC$ be an equilateral triangle with side length $55.$ Points $D,$ $E,$ and $F$ lie on $\overline{BC},$ $\overline{CA},$ and $\overline{AB},$ respectively, with $BD = 7,$ $CE=30,$ and $AF=40.$ Point $P$ inside $\triangle ABC$ has the property that 
 \[\angle AEP = \angle BFP = \angle CDP.\] Find $\tan^2(\angle AEP).$


 
 \textbf{Answer: }075



\textbf{Solution 1 (Coordinates Bash)}

By Miquel's theorem, $P=(AEF)\cap(BFD)\cap(CDE)$ (intersection of circles)$\text{*}$. The law of cosines can be used to compute $DE=42$, $EF=35$, and $FD=13$. Toss the points on the coordinate plane; let $B=(-7, 0)$, $D=(0, 0)$, and $C=(48, 0)$, where we will find $\tan^{2}\left(\measuredangle CDP\right)$ with $P=(BFD)\cap(CDE)$.

 By the extended law of sines, the radius of circle $(BFD)$ is $\frac{13}{2\sin 60^{\circ}}=\frac{13}{3}\sqrt{3}$. Its center lies on the line $x=-\frac{7}{2}$, and the origin is a point on it, so $y=\frac{23}{6}\sqrt{3}$.

 The radius of circle $(CDE)$ is $\frac{42}{2\sin 60^{\circ}}=14\sqrt{3}$. The origin is also a point on it, and its center is on the line $x=24$, so $y=2\sqrt{3}$.

 The equations of the two circles are 
 \begin{align*}(BFD)&:\left(x+\frac{7}{2}\right)^{2}+\left(y-\frac{23}{6}\sqrt{3}\right)^{2}=\frac{169}{3} \\ (CDE)&:\left(x-24\right)^{2}+\left(y-2\sqrt{3}\right)^{2}=588\end{align*}
 These equations simplify to 
 \begin{align*}(BFD)&:x^{2}+7x+y^{2}-\frac{23}{3}\sqrt{3}y=0 \\ (CDE)&: x^{2}-48x+y^{2}-4\sqrt{3}y=0\end{align*}
 Subtracting these two equations gives that both their points of intersection, $D$ and $P$, lie on the line $55x-\frac{11}{3}\sqrt{3}y=0$. Hence, $\tan^{2}\left(\measuredangle AEP\right)=\tan^{2}\left(\measuredangle CDP\right)=\left(\frac{55}{\frac{11}{3}\sqrt{3}}\right)^{2}=3\left(\frac{55}{11}\right)^{2}=\boxed{075}$. To scale, the configuration looks like the figure below:
                    \begin{figure}[H]
                        \centering
                        \includegraphics[width=0.5\textwidth]{Solution/AIME_I_2023_12_Solution_Figure_1.png}
                    \end{figure}
                    

 $\text{*}$Basic angle chasing gives $\angle BDP=180-\angle CDP.$ Because $\angle BFP+\angle BDP=180+\angle BFP-\angle CDP=180,$ which means that $BFPD$ is cyclic, and that $P$ passes through the circumcircle of triangle $BFD.$ Similar reasoning leads us to the fact that $P$ also passes through the circumcircles of triangles $AEF$ and $CDE,$ which means that $P=(AEF)\cap(BFD)\cap(CDE).$ Continue as above.

 Note: Since $D$ and $P$ are the two intersections of $(BFD)$ and $(CDE),$ we know that $DP$ is perpendicular to the line between the centers of $(BFD)$ and $(CDE).$ Thus, after getting the coordinates of the centers of $(BFD)$ and $(CDE),$ we can then immediately find the slope of line $DP,$ and then get the answer from there.

 



\textbf{Solution 2 (Vectors/Complex)}

Denote $\theta = \angle AEP$.

 In $AFPE$, we have $\overrightarrow{AF} + \overrightarrow{FP} + \overrightarrow{PE} + \overrightarrow{EA} = 0$.Thus,
 \[ AF + FP e^{i \theta} + PE e^{i \left( \theta + 60^\circ \right)} + EA e^{- i 120^\circ} = 0. \]

 Taking the real and imaginary parts, we get
 \begin{align*} AF + FP \cos \theta + PE \cos \left( \theta + 60^\circ \right) + EA \cos \left( - 120^\circ \right) & = 0 \hspace{1cm} (1) \\ FP \sin \theta + PE \sin \left( \theta + 60^\circ \right) + EA \sin \left( - 120^\circ \right) & = 0 \hspace{1cm} (2) \end{align*}


 In $BDPF$, analogous to the analysis of $AFPE$ above, we get
 \begin{align*} BD + DP \cos \theta + PF \cos \left( \theta + 60^\circ \right) + FB \cos \left( - 120^\circ \right) & = 0 \hspace{1cm} (3) \\ DP \sin \theta + PF \sin \left( \theta + 60^\circ \right) + FB \sin \left( - 120^\circ \right) & = 0 \hspace{1cm} (4) \end{align*}


 Taking $(1) \cdot \sin \left( \theta + 60^\circ \right) - (2) \cdot \cos \left( \theta + 60^\circ \right)$, we get
 \[ AF \sin \left( \theta + 60^\circ \right) + \frac{\sqrt{3}}{2} FP - EA \sin \theta = 0 . \hspace{1cm} (5) \]

 Taking $(3) \cdot \sin \theta - (4) \cdot \cos \theta$, we get
 \[ BD \sin \theta - \frac{\sqrt{3}}{2} FP + FB \sin \left( \theta + 120^\circ \right) . \hspace{1cm} (6) \]

 Taking $(5) + (6)$, we get
 \[ AF \sin \left( \theta + 60^\circ \right)  - EA \sin \theta + BD \sin \theta  + FB \sin \left( \theta + 120^\circ \right) . \]

 Therefore,
 \begin{align*} \tan \theta & = \frac{\frac{\sqrt{3}}{2} \left( AF + FB \right)} {\frac{FB}{2} + EA - \frac{AF}{2} - BD} \\ & = 5 \sqrt{3} . \end{align*}


 Therefore, $\tan^2 \theta = \boxed{075}$.

 



\textbf{Solution 3 (Synthetic)}

Drop the perpendiculars from $P$ to $\overline{AB}$, $\overline{AC}$, $\overline{BC}$, and call them $Q,R,$ and $S$ respectively. This gives us three similar right triangles $FQP$, $ERP$, and $DSP.$ 

 The sum of the perpendiculars to a point $P$ within an equilateral triangle is always constant, so we have that $PQ+PR+PS=\dfrac{55 \sqrt{3}}{2}.$ 

 The sum of the lengths of the alternating segments split by the perpendiculars from a point $P$ within an equilateral triangle is always equal to half the perimeter, so $QA+RC+SB = \dfrac{165}{2},$ which means that $FQ+ER+DS = QA+RC+SB - CE - AF - BD = \dfrac{165}{2} - 30 - 40 - 7 = \dfrac{11}{2}.$

 Finally, $\tan AEP = \dfrac{PQ}{FQ} = \dfrac{PR}{ER} = \dfrac{PS}{DS} = \dfrac{PQ+PR+PS}{FQ+ER+DS} = 5 \sqrt{3}.$ 

 Thus, $\tan^2 AEP = \boxed{075}.$

 <i>\textbf{Claim}</i>

 
                    \begin{figure}[H]
                        \centering
                        \includegraphics[width=0.5\textwidth]{Solution/AIME_I_2023_12_Solution_Figure_2.png}
                    \end{figure}
                    a) Carnot's theorem. Given triangle $\triangle ABC$ and point $P.$  Let $PS \perp BC,$$PR \perp AC, PQ \perp AB.$ $P$ doesn't have to be inside $\triangle ABC.$

 Prove that $AQ^2 + BS^2 + CR^2 = AR^2 + BQ^2 + CS^2.$

 b) Let $\triangle ABC$ be the equilateral triangle. Prove that $AQ + BS + CR = \frac {3}{2} AB.$ (The sum of the lengths of the alternating segments split by the perpendiculars from a point $P$ within an equilateral triangle is equal to half the perimeter.)

 <i>\textbf{Proof}</i>

 a) $AP^2 = AQ^2 + PQ^2 = AR^2 + PR^2,$
 \[BP^2 = BQ^2 + PQ^2 = BS^2 + PS^2,\]
 \[CP^2 = CR^2 + PR^2 = CS^2 + PS^2 \implies\]$AQ^2 + PQ^2 + BS^2 + PS^2 + CR^2 + PR^2 = AR^2 + PR^2 + BQ^2 + PQ^2 + CS^2 + PS^2,$
 \[AQ^2 + BS^2 + CR^2 = AR^2 + BQ^2 + CS^2.\]b) $AQ + BQ = BS + CS = CR + AR = AB = AC = BC,$
 \[AQ^2 + BQ^2 + 2 AQ \cdot BQ + BS^2 + CS^2 + 2 BS \cdot CS + AR^2 + CR^2+ 2 AR \cdot CR = 3 AB^2.\]

 
                    \begin{figure}[H]
                        \centering
                        \includegraphics[width=0.5\textwidth]{Solution/AIME_I_2023_12_Solution_Figure_3.png}
                    \end{figure}
                    
 \[2AQ^2 + 2 AQ \cdot BQ + 2 BS^2 + 2 BS \cdot CS + 2 CR^2 + 2 AR \cdot CR = 3 AB^2.\]
 \[AQ (AQ + BQ) + BS (BS + CS) + CR (AR + CR) = \frac {3}{2} AB^2.\]
 \[AQ \cdot AB + BS \cdot BC + CR \cdot AC = \frac {3}{2} AB^2.\]
 \[AQ + BS + CR = \frac {3}{2} AB.\]

 \textbf{vladimir.shelomovskii@gmail.com, vvsss}

 



\textbf{Solution 4 (Law of Cosines)}

This solution is inspired by AIME 1999 Problem 14 Solution 2 (a similar question)

 Draw line segments from $P$ to points $A$, $B$, and $C$. And label the angle measure of $\angle{BFP}$, $\angle{CDP}$, and $\angle{AEP}$ to be $\alpha$

 Using Law of Cosines (note that $\cos{\angle{AFP}}=\cos{\angle{BDP}}=\cos{\angle{CEP}}=\cos{180^\circ-\alpha}=-\cos{\alpha}$)

 
 \begin{align*} (1) \ BP^2 &= FP^2+15^2-2\cdot FP\cdot15\cdot\cos(\alpha)\\ (2) \ BP^2 &= DP^2+7^2+2\cdot DP\cdot7\cdot\cos(\alpha)\\ (3) \ CP^2 &= DP^2+48^2-2\cdot DP\cdot48\cdot\cos(\alpha)\\ (4) \ CP^2 &= EP^2+30^2+2\cdot EP\cdot30\cdot\cos(\alpha)\\ (5) \ AP^2 &= EP^2+25^2-2\cdot EP\cdot25\cdot\cos(\alpha)\\ (6) \ AP^2 &= FP^2+40^2+2\cdot FP\cdot40\cdot\cos(\alpha)\\ \end{align*}


 We can perform this operation $(1) - (2) + (3) - (4) + (5) - (6)$:

 Leaving us with (after combining and simplifying)
 \[\cos{\alpha}=\frac{11}{2\cdot(DP+EP+FP)}\]

 Therefore, we want to solve for $DP+EP+FP$

 Notice that $\angle{DPE}=\angle{EPF}=\angle{FPD}=120^\circ$

 We can use Law of Cosines again to solve for the sides of $\triangle{DEF}$, which have side lengths of $13$, $42$, and $35$, and area $120\sqrt{3}$.

 Label the lengths of $PD$, $PE$, and $PF$ to be $x$, $y$, and $z$.

 Therefore, using the $\sin$ area formula,

 
 \begin{align*} [\triangle{DEF}] &= \frac{1}{2}\cdot\sin{120°}\cdot(xy+yz+zx) = 120\sqrt{3} \\ xy+yz+zx &= 2^5\cdot3\cdot5 \end{align*}


 In addition, we know that

 
 \begin{align*} x^2+y^2+xy&=42^2\\ y^2+z^2+yz&=35^2\\ z^2+x^2+zx&=13^2\\ \end{align*}


 By using Law of Cosines for $\triangle{DPE}$, $\triangle{EPF}$, and $\triangle{FPD}$ respectively

 Because we want $DP+EP+FP$, which is $x+y+z$, we see that

 
 \begin{align*} (x+y+z)^2 &= \frac{(x^2+y^2+xy)+(y^2+z^2+yz)+(z^2+x^2+zx)+3(xy+yz+zx)}{2} \\ (x+y+z)^2 &= \frac{42^2+35^2+13^2+3\cdot2^5\cdot3\cdot5}{2} \\ (x+y+z)^2 &= 2299 \\ x+y+z &= 11\sqrt{19} \end{align*}


 So plugging the results back into the equation before, we get

 
 \begin{align*} \cos{\alpha} &= \frac{-1}{2\sqrt{19}}\\ \sin{\alpha} &= \frac{5\sqrt{3}}{2\sqrt{19}} \end{align*}


 Giving us 
 \[\tan^2{\alpha}=\boxed{075}\]

 



\textbf{Solution 5 (Combining Solutions 3 and 4)}

We begin by using the fact stated in Solution 3 that, for any point in an equilateral triangle, the lengths of the three perpendicular lines dropped to the sides of the triangle add up to the altitude of that triangle. To make things simple, let's assign $\angle AEP = \angle BFP = \angle CDP = \alpha$. We can label these three perpendiculars as:
 \[PD\cdot\sin{\alpha} + PE\cdot\sin{\alpha} + PF\cdot\sin{\alpha} = \dfrac{55 \sqrt{3}}{2}\]Simplifying, we get
 \[PD + PE + PF = \dfrac{55 \sqrt{3}}{2\cdot\sin{\alpha}}\]Now, as stated and quoting Solution 4,"Draw line segments from $P$ to points $A$, $B$, and $C$. [We know that] the angle measure of $\angle{AEP}$, $\angle{BFP}$, and $\angle{CDP}$ is $\alpha$

 Using Law of Cosines (note that $\cos{\angle{AFP}}=\cos{\angle{BDP}}=\cos{\angle{CEP}}=\cos{180^\circ-\alpha}=-\cos{\alpha}$)

 
 \begin{align*} (1) \ BP^2 &= FP^2+15^2-2\cdot FP\cdot15\cdot\cos(\alpha)\\ (2) \ BP^2 &= DP^2+7^2+2\cdot DP\cdot7\cdot\cos(\alpha)\\ (3) \ CP^2 &= DP^2+48^2-2\cdot DP\cdot48\cdot\cos(\alpha)\\ (4) \ CP^2 &= EP^2+30^2+2\cdot EP\cdot30\cdot\cos(\alpha)\\ (5) \ AP^2 &= EP^2+25^2-2\cdot EP\cdot25\cdot\cos(\alpha)\\ (6) \ AP^2 &= FP^2+40^2+2\cdot FP\cdot40\cdot\cos(\alpha)\\ \end{align*}


 We can perform this operation $(1) - (2) + (3) - (4) + (5) - (6)$:

 Leaving us with (after combining and simplifying)
 \[PD + PE + PF=\frac{11}{2\cdot\cos{\alpha}}\]".

 Now, we can use our previous equation along with this one to get:
 \[\frac{11}{2\cdot\cos{\alpha}} = \dfrac{55 \sqrt{3}}{2\cdot\sin{\alpha}}\].

 This equation becomes:
 \[\tan{\alpha} = 5\sqrt{3}\]As so, our answer is
 \[\left(5\sqrt3\right)^2=\boxed{075}.\]~Solution by armang32324 (Mathemagics Club)

 



\textbf{Solution 6}

By the law of cosines,
 \[FE=\sqrt{AF^2+AE^2-2AF\cdot AE\cos\angle FAE}=35.\]Similarly we get $FD=13$ and $DE=42$. $\angle AEP=\angle BFP=\angle CDP\overset\triangle=\theta$ implies that $AFPE$, $BDPF$, and $CDPE$ are three cyclic quadrilaterals, as shown below:
                    \begin{figure}[H]
                        \centering
                        \includegraphics[width=0.5\textwidth]{Solution/AIME_I_2023_12_Solution_Figure_4.png}
                    \end{figure}
                    Using the law of sines in each,
 \[\frac{AP}{35}=\frac{AP}{FE}=\frac{CP}{42}=\frac{CP}{ED}=\frac{BP}{13}=\frac{BP}{DF}=\frac{\sin\theta}{\sin\frac\pi3}.\]So we can set $AP=35k$, $BP=13k$, and $CP=42k$. Let $PD=d$, $PE=e$, and $PF=f$. Applying Ptolemy theorem in the cyclic quadrilaterals,
 \[\begin{cases}AP\cdot FE=AF\cdot PE+AE\cdot PF,\\CP\cdot ED=CE\cdot PD+CD\cdot PE,\\BP\cdot DF=BD\cdot PF+BF\cdot PD.\end{cases} \implies \begin{cases} 1225k=40e+25f,\\1764k=30d+48e,\\169k=15d+7f, \end{cases}\]We can solve out $d=\frac{54k}5$, $e=30k$, $f=k$. By the law of cosines in $\triangle PEF$, $FE=\sqrt{900k^2+k^2-60k\cdot\left(\frac{-1}2\right)}=\sqrt{931}k$. The law of sines yield $\frac{\sin\angle AEP}{\sin\angle FAE}=\frac{AP}{FE}=\frac{35k}{\sqrt{931}k}=\frac{35}{\sqrt{931}}$.Lastly, $\sin\angle AEP=\frac{5\sqrt{57}}{38}$, then $\tan\angle AEP=5\sqrt3$. The answer is
 \[\left(5\sqrt3\right)^2=\boxed{075}.\]

 



\textbf{Solution 7}

$P$ is well known to be the Miquel point of triangle $DEF$. By Law of Cosines, $E F = 35$, $F D = 13$, $D E = 42$. Call $P D$, $P E$, $P F$: $x$, $y$, $z$ respectively. Then we obtain that $x^2+xy+y^2 = 42^2$, $x^2+xz+z^2 = 13^2$, $y^2+yz+z^2 = 35^2$. Note that since the area of the triangle by sine area formula is $(xy + xz + yz)\left(\frac{\sqrt{3}}{4}\right)$, but by Heron's formula, the area of the triangle is $\sqrt{(45)(3)(32)(10)} = (5)(3)(8)\sqrt{3} = 120\sqrt{3}$. $xy+xz+yz = 480$, and so we obtain $x+y+z = 11\sqrt{19}$ by adding all the equations together. Now, subtract any two equations from each other to obtain $(y-z)(11\sqrt{19}) = (55)(29)$, $(x-z)(11\sqrt{19}) = (7)(77)$, $x-z = \frac{49}{\sqrt{19}}$, $y-z = \frac{145}{\sqrt{19}}$, and $z = \frac{5}{\sqrt{19}}$, $y = \frac{150}{\sqrt{19}}$, and $x = \sqrt{54}{\sqrt{19}}$. Now we find PC through Ptolemy's theorem, do Law of Cosines, and we are finished.

 



\textbf{Solution 8 (Coordinates Bash Again!)}

Recall that if $\theta$ is the angle between two lines with slopes $m_1, m_2$, we can write:

 $\tan(\theta) = \mid \frac{m_1 - m_2}{1 + m_1 \cdot m_2} \mid$ 

 We will use this fact to solve the problem. Let $B$ be at the origin, $C$ be at $(55, 0)$ and $A$ be at $(\frac{55}{2}, \frac{55\sqrt{3}}{2})$. Next, we can easily find the equation of line $AB$ as $y = x\sqrt{3}$ and the equation of line $AC$ as $y = -x\sqrt{3} + 55\sqrt{3}$. We also know that $D$ is at $(7, 0)$ given the problem's conditions. We can now find the coordinates of $E$ and $F$. To find the coordinates of $F$, recall that $F$ lies on line segment $AB$ which has equation $y = x\sqrt{3}$. Hence, $F$ will be a point of the form $(x, x\sqrt{3})$. We also know the distance between $F$ and $B$(the origin), is just $55 - 40 = 15$. So the distance between the point $(x, x\sqrt{3})$ and $(0, 0)$ is just $15$. We can use the distance formula to get a solution of $x = \frac{15}{2}$(there will be two solutions but one of them is unfeasible). So the coordinates of $F$ is just $(\frac{15}{2}, \frac{15\sqrt{3}}{2})$. Similarly, we can find the coordinates of $E$ in a similar way to get $(40, 15\sqrt{3})$. Note that $\angle BFP = \angle PEA = \angle PDC = \theta$. Let's focus on $\angle BFP$. This angle is between lines $BF$ and $FP$ so we just have to compute the slopes of these lines. Let $P$ be at point $(a, b)$. Then, the slope of $BF$ is just $\sqrt{3}$ and the slope of $FP$ is just $\frac{2b - 15\sqrt{3}}{2a - 15}$. Because the angle between these two lines is $\theta$, we can use our above lemma to get the equation:

 $\tan(\theta) = \mid \frac{\sqrt{3} - \frac{2b - 15\sqrt{3}}{2a - 15}}{1 + \sqrt{3} \cdot \frac{2b - 15\sqrt{3}}{2a - 15}} \mid$

 We can similarly use the other angles and lines to come up with two more of these type of equations:

 $\tan(\theta) = \mid \frac{-\sqrt{3} - \frac{15\sqrt{3} - b}{40 - a}}{1 - \sqrt{3} \cdot \frac{15\sqrt{3} - b}{40 - a}} \mid$

 $\tan(\theta) = \mid \frac{b}{7 - a} \mid$

 We have $3$ equations with $3$ variables and since we are trying to find $\tan(\theta)^{2}$, we can treat $\tan(\theta)$ as a variable. We solve to get $\tan(\theta) = 5\sqrt{3}$ and therefore our answer is just $\boxed{075}$.

 $\textbf{Note: How to solve the system of equations}$

 The first step is to realize that we need to make the inputs nicer. So we clear the fractions and set them equal. Doing so, we have:

 $\frac{a\sqrt{3} - b}{a + b\sqrt{3} - 30} = \frac{a\sqrt{3} + b - 55\sqrt{3}}{-a + b\sqrt{3} - 5} = \frac{b}{7 - a}$. 

 Now we can write:

 $\frac{a\sqrt{3} - b}{a + b\sqrt{3} - 30} = \frac{b}{7 - a}$

 $\frac{a\sqrt{3} + b - 55\sqrt{3}}{-a + b\sqrt{3} - 5} = \frac{b}{7 - a}$

 Here, we're going to change $a$ to $x$ and $b$ to $y$ just to make the solving more traditional. We can cross multiply the first equation and simplify to get 

 $7x\sqrt{3} - x^{2} \sqrt{3} = y^{2} \sqrt{3} - 23y$

 Next, we cross multiply the second equation to get:

 $7x\sqrt{3} - x^{2} \sqrt{3} + 7y - xy - 385\sqrt{3} + 55x\sqrt{3} = -xy + y^{2} \sqrt{3} - 5y$

 Notice we could make some substitutions from our first equation to the second equation:

 $y^{2} \sqrt{3} - 23y + 7y - 385\sqrt{3} + 55x\sqrt{3} = y^{2} \sqrt{3} - 5y$

 We simplify this to get:

 $-16y - 385\sqrt{3} + 55x\sqrt{3} = -5y$

 This means:

 $55x\sqrt{3} - 385\sqrt{3} = 11y$

 Which yields:

 $y = 5\sqrt{3}(x - 7)$

 Now, we plug this representation into the equation $7x\sqrt{3} - x^{2} \sqrt{3} = y^{2} \sqrt{3} - 23y$ we got earlier to get:

 $7x\sqrt{3} - x^{2} \sqrt{3} = 75\sqrt{3} (x - 7)^{2} - 115x\sqrt{3} + 805\sqrt{3}$

 We can expand and simplify to get:

 $122x\sqrt{3} - x^{2} \sqrt{3} = 75x^{2} \sqrt{3} - 1050x\sqrt{3} + 4480\sqrt{3}$

 Simplifying, we get:

 $76x^{2} \sqrt{3} - 1172x\sqrt{3} + 4480\sqrt{3} = 0$

 This yields:

 $19x^{2} - 293x + 1120 = 0$

 Now this is trivial to solve:

 $x = \frac{293 \pm \sqrt{293^{2} - 76\cdot1120}}{38}$

 To find $293^{2}$ really quickly under contest conditions, we can do:

 $293^{2} = (200 + 93)^{2} = 200^{2} + 400\cdot93 + (90 + 3)^{2} = 40000 + 37200 + 8100 + 549 = 85849$

 Now back to simplifying $x$:

 The discriminant nicely simplifies to $85849 - 85120 = 729 = 27^{2}$

 So we have $x = \frac{293 + 27}{38} = \frac{320}{38} = \frac{160}{19}$. (Note that if we take the negative version of the solution, we would have $x = \frac{266}{38} = 7$ but this would yield $y = 5\sqrt{3} \cdot 7 - 35\sqrt{3} = 0$ which is absurd)

 And now we're basically done. Recall that $y = 5x\sqrt{3} - 35\sqrt{3} = \frac{5\sqrt{3}\cdot160}{19} - \frac{35\sqrt{3}\cdot19}{19} = \frac{135\sqrt{3}}{19}$. Now recall that we let $a = x$ and $b = y$ so we now have $a = \frac{160}{19}$ and $b = \frac{135\sqrt{3}}{19}$. 

 Now, recall that $\tan(\theta) = \mid \frac{b}{7 - a} \mid = \mid \frac{\frac{135\sqrt{3}}{19}}{7 - \frac{160}{19}} \mid = \mid \frac{\frac{135\sqrt{3}}{19}}{\frac{133 - 160}{19}} \mid = \mid \frac{135\sqrt{3}}{-27} \mid = \mid -5\sqrt{3} \mid = 5\sqrt{3}$ as desired.

 


	\item (\textit{AIME I 2023}) Each face of two noncongruent parallelepipeds is a rhombus whose diagonals have lengths $\sqrt{21}$ and $\sqrt{31}$. The ratio of the volume of the larger of the two polyhedra to the volume of the smaller is $\frac{m}{n}$, where $m$ and $n$ are relatively prime positive integers. Find $m + n$. A parallelepiped is a solid with six parallelogram faces such as the one shown below.


 
                    \begin{figure}[H]
                        \centering
                        \includegraphics[width=0.5\textwidth]{Problem/AIME_I_2023_13_Problem_Figure_1.png}
                    \end{figure}
                    


 
 \textbf{Answer: }125



\textbf{Solution 1 (3-D Vector Analysis)}

Denote $\alpha = \tan^{-1} \frac{\sqrt{21}}{\sqrt{31}}$.Denote by $d$ the length of each side of a rhombus.

 Now, we put the solid to the 3-d coordinate space.We put the bottom face on the $x-O-y$ plane.For this bottom face, we put a vertex with an acute angle $2 \alpha$ at the origin, denoted as $O$.For two edges that are on the bottom face and meet at $O$, we put one edge on the positive side of the $x$-axis. The endpoint is denoted as $A$. Hence, $A = \left( d , 0 , 0 \right)$.We put the other edge in the first quadrant of the $x-O-y$ plane. The endpoint is denoted as $B$. Hence, $B = \left( d \cos 2 \alpha , d \sin 2 \alpha , 0 \right)$.

 For the third edge that has one endpoint $O$, we denote by $C$ its second endpoint.We denote $C = \left( u , v , w \right)$.Without loss of generality, we set $w > 0$.Hence,
 \[ u^2 + v^2 + w^2 = d^2 . \hspace{1cm} (1) \]

 We have
 \begin{align*} \cos \angle AOC & = \frac{\overrightarrow{OA} \cdot \overrightarrow{OC}}{|OA| \cdot |OC|} \\ & = \frac{u}{d} , \hspace{1cm} (2) \end{align*}
and
 \begin{align*} \cos \angle BOC & = \frac{\overrightarrow{OB} \cdot \overrightarrow{OC}}{|OB| \cdot |OC|} \\ & = \frac{u \cos 2 \alpha + v \sin 2 \alpha}{d} . \hspace{1cm} (3) \end{align*}


 Case 1: $\angle AOC = \angle BOC = 2 \alpha$ or $2 \left( 90^\circ - \alpha \right)$.

 By solving (2) and (3), we get
 \begin{align*} u & = \pm d \cos 2 \alpha , \\ v & = \pm d \cos 2 \alpha \frac{1 - \cos 2 \alpha}{\sin 2 \alpha} \\ & = \pm d \cos 2 \alpha \tan \alpha . \end{align*}


 Plugging these into (1), we get
 \begin{align*} w & = d \sqrt{1 - \cos^2 2 \alpha - \cos^2 2 \alpha \tan^2 \alpha} \\ & = d \sqrt{\sin^2 2 \alpha - \cos^2 2 \alpha \tan^2 \alpha} . \hspace{1cm} (4) \end{align*}


 Case 2: $\angle AOC = 2 \alpha$ and $\angle BOC = 2 \left( 90^\circ - \alpha \right)$, or $\angle BOC = 2 \alpha$ and $\angle AOC = 2 \left( 90^\circ - \alpha \right)$.

 By solving (2) and (3), we get
 \begin{align*} u & = \pm d \cos 2 \alpha , \\ v & = \mp d \cos 2 \alpha \frac{1 + \cos 2 \alpha}{\sin 2 \alpha} \\ & = \mp d \cos 2 \alpha \cot \alpha . \end{align*}


 Plugging these into (1), we get
 \begin{align*} w & = d \sqrt{1 - \cos^2 2 \alpha - \cos^2 2 \alpha \cot^2 \alpha} \\ & = d \sqrt{\sin^2 2 \alpha - \cos^2 2 \alpha \cot^2 \alpha} . \hspace{1cm} (5) \end{align*}


 We notice that $(4) > (5)$. Thus, (4) (resp. (5)) is the parallelepiped with a larger (resp. smaller) height.

 Therefore, the ratio of the volume of the larger parallelepiped to the smaller one is
 \begin{align*} \frac{(4)}{(5)} & = \frac{\sqrt{\sin^2 2 \alpha - \cos^2 2 \alpha \tan^2 \alpha}} {\sqrt{\sin^2 2 \alpha - \cos^2 2 \alpha \cot^2 \alpha}} \\ & = \sqrt{\frac{\tan^2 2 \alpha - \tan^2 \alpha}{\tan^2 2 \alpha - \cot^2 \alpha}} . \end{align*}


 Recall that $\tan \alpha = \frac{\sqrt{21}}{\sqrt{31}}$.Thus, $\tan 2 \alpha = \frac{2 \tan \alpha}{1 - \tan^2 \alpha} = \frac{\sqrt{21 \cdot 31}}{5}$.Plugging this into the equation above, we get
 \begin{align*} \frac{(4)}{(5)} & = \frac{63}{62}. \end{align*}


 Therefore, the answer is $63 + 62 = \boxed{\textbf{(125) }}$.

 



\textbf{Solution 2 (no trig)}

Let one of the vertices be at the origin and the three adjacent vertices be $u$, $v$, and $w$. For one of the parallelepipeds, the three diagonals involving the origin have length $\sqrt {21}$. Hence, $(u+v)\cdot (u+v)=u\cdot u+v\cdot v+2u\cdot v=21$ and $(u-v)\cdot (u-v)=u\cdot u+v\cdot v-2u\cdot v=31$. Since all of $u$, $v$, and $w$ have equal length, $u\cdot u=13$, $v\cdot v=13$, and $u\cdot v=-2.5$. Symmetrically, $w\cdot w=13$, $u\cdot w=-2.5$, and $v\cdot w=-2.5$. Hence the volume of the parallelepiped is given by $\sqrt{\operatorname{det}\begin{pmatrix}13&-2.5&-2.5\\-2.5&13&-2.5\\-2.5&-2.5&13\end{pmatrix}}=\sqrt{\operatorname{det}\begin{pmatrix}15.5&-15.5&0\\-2.5&13&-2.5\\0&-15.5&15.5\end{pmatrix}}=\sqrt{15.5^2\operatorname\det\begin{pmatrix}1&-1&0\\-2.5&13&-2.5\\0&-1&1\end{pmatrix}}=\sqrt{15.5^2\cdot 8}$.

 For the other parallelepiped, the three diagonals involving the origin are of length $\sqrt{31}$ and the volume is $\sqrt{\operatorname{det}\begin{pmatrix}13&2.5&2.5\\2.5&13&2.5\\2.5&2.5&13\end{pmatrix}}=\sqrt{\operatorname{det}\begin{pmatrix}10.5&-10.5&0\\2.5&13&2.5\\0&-10.5&10.5\end{pmatrix}}=\sqrt{10.5^2\operatorname\det\begin{pmatrix}1&-1&0\\2.5&13&2.5\\0&-1&1\end{pmatrix}}=\sqrt{10.5^2\cdot 18}$.

 Consequently, the answer is $\sqrt\frac{10.5^2\cdot 18}{15.5^2\cdot 8}=\frac{63}{62}$, giving $\boxed{125}$.

 Note on the linear algebra formula:It is well known that a parallelepiped whose 3 adjacent vertices have vectors $u= \langle u_1, u_2, u_3 \rangle$, $v= \langle v_1, v_2, v_3 \rangle$, and $w = \langle w_1, w_2, w_3 \rangle$ volume is the absolute value of$\operatorname{det} \begin{pmatrix} u_1 & u_2 & u_3 \\ v_1 & v_2 & v_3 \\ w_1 & w_2 & w_3  \end{pmatrix}$Then the volume is also

 \begin{align*}\sqrt{\left (\operatorname{det} \begin{pmatrix}u_1 & u_2 & u_3 \\v_1 & v_2 & v_3 \\w_1 & w_2 & w_3 \\\end{pmatrix}\right)^2} &= \sqrt{\operatorname{det} \left( \begin{pmatrix}u_1 & u_2 & u_3 \\v_1 & v_2 & v_3 \\w_1 & w_2 & w_3 \\\end{pmatrix} \cdot \begin{pmatrix}u_1 & u_2 & u_3 \\v_1 & v_2 & v_3 \\w_1 & w_2 & w_3 \\\end{pmatrix} \right) } \\&= \sqrt{\operatorname{det} \left( \begin{pmatrix}u_1 & u_2 & u_3 \\v_1 & v_2 & v_3 \\w_1 & w_2 & w_3 \\\end{pmatrix} \cdot \begin{pmatrix}u_1 & v_1 & w_1 \\u_2 & v_2 & w_2 \\u_3 & v_3 & w_3 \\\end{pmatrix} \right) } = \sqrt{\operatorname{det} \begin{pmatrix}u \cdot u & u \cdot v & u \cdot w \\v \cdot u & v \cdot v & v \cdot w \\w \cdot u & w \cdot v & w \cdot w\end{pmatrix} } \end{align*}


 Since the determinant of the transpose of a matrix equals to the determinant of the original matrix.

 



\textbf{Solution 3 (No trig, no linear algebra)}

Observe that both parallelepipeds have two vertices (one on each base) that have three congruent angles meeting at them. Denote the parallelepiped with three acute angles meeting $P$, and the one with three obtuse angles meeting $P'$.

 The volume of a parallelepiped is simply the base area times the height, but because both parallelepipeds have the same base, what we want is just the ratio of the heights.

 Denote the point with three acute angles meeting at it in $P$ as $A$, and its neighbors $B$, $C$, and $D$. Similarly, denote the point with three obtuse angles meeting at it in $P'$ as $A'$, and its neighbors $B'$, $C'$, and $D'$.

 We have the following equations:

 
 \[\textrm{Height of }P\textrm{ from }ACD = \frac{\textrm{Vol}(ABCD) \cdot 3}{[ACD]},\]
 \[\textrm{Height of }P'\textrm{ from }A'C'D' = \frac{\textrm{Vol}(A'B'C'D') \cdot 3}{[A'C'D']}.\]

 However, $ACD$ and $A'C'D'$ are both half the area of a rhombus with diagonals $\sqrt{31}$ and $\sqrt{21}$, so our ratio is really

 
 \[\frac{P}{P'} = \frac{\textrm{Vol}(ABCD)}{\textrm{Vol}(A'B'C'D')}.\]

 Because the diagonals of all of the faces are $\sqrt{31}$ and $\sqrt{21}$, each edge of the parallelepipeds is $\sqrt{13}$ by the Pythagorean theorem. 

 We have $AB = AC = AD = \sqrt{13}$, and $BC = CD = BD = \sqrt{21}$. When we drop a perpendicular to the centroid of $BCD$ from $A$ (let's call this point $O$), we have $BO = \frac{\sqrt{21}}{\sqrt{3}} = \sqrt{7}$. Thus, 

 
 \[AB^2 - BO^2 = AO^2\]
 \[13 - 7 = AO^2 = 6\]
 \[AO = \sqrt{6}.\]

 The area of base $BCD$ is $\frac{21\sqrt{3}}{4}$. Hence, 

 
 \[\textrm{Vol}(ABCD) = \frac{\sqrt{6}\cdot\frac{21\sqrt{3}}{4}}{3}\]
 \[= \frac{63\sqrt{2}}{12}.\]

 We can apply a similar approach to $A'B'C'D'$.

 $A'B' = A'C' = A'D' = \sqrt{13}$, and $B'C' = C'D' = B'D' = \sqrt{31}$. When we drop a perpendicular to the centroid of $B'C'D'$ from $A'$ (let's call this point $O'$), we have $B'O' = \frac{\sqrt{31}}{\sqrt{3}} = \sqrt{\frac{31}{3}}$. Thus, 

 
 \[A'B'^2 - B'O'^2 = A'O'^2\]
 \[13 - \frac{31}{3} = A'O'^2\]
 \[A'O' = \sqrt{\frac{8}{3}} = \frac{2\sqrt{6}}{3}.\]

 The area of base $B'C'D'$ is $\frac{31\sqrt{3}}{4}$. Hence, 

 
 \[\textrm{Vol}(A'B'C'D') = \frac{\frac{2\sqrt{6}}{3}\cdot\frac{31\sqrt{3}}{4}}{3}\]
 \[= \frac{186\sqrt{2}}{36}\]
 \[= \frac{62\sqrt{2}}{12}.\]

 Finally, 

 
 \[\frac{P}{P'} = \frac{\textrm{Vol}(ABCD)}{\textrm{Vol}(A'B'C'D')} = \frac{\frac{63\sqrt{2}}{12}}{\frac{62\sqrt{2}}{12}} = \frac{63}{62}.\]

 Our answer is $63 + 62 = \boxed{125}$.

 



\textbf{Solution 4 (Pythagorean theorem)}

Since the two parallelepipeds have the same base, all we need to do is to find their respective heights.

 
                    \begin{figure}[H]
                        \centering
                        \includegraphics[width=0.5\textwidth]{Solution/AIME_I_2023_13_Solution_Figure_1.png}
                    \end{figure}
                    

 As illustrated in the above diagram, drop a perpendicular from $D$ onto the base at $P$. Denote the center of the base by $O$. By symmetry, $P$ must be on $AO$. Now we need to find $DP$.

 Apply Pythagorean theorem to $\triangle DPA$ we have
 \[DP^2 = AD^2 - AP^2.\]

 Apply Pythagorean theorem to $\triangle DPC$ and then $\triangle CPO$ we have
 \[DP^2 = DC^2 - CP^2 = DC^2 - (CO^2 + OP^2) = DC^2 - (CO^2 + (AO-AP)^2) = DC^2 - CO^2 - (AO-AP)^2.\]

 Combining the above two, we have
 \[AD^2 - AP^2 = DC^2 - CO^2 - (AO-AP)^2.\]

 Since $AD=\sqrt{13}$, $DC=\sqrt{21}$, $CO=\frac{\sqrt{21}}{2}$, $AO=\frac{\sqrt{31}}{2}$, plug them into the above equation and solve for the only unknown variable $AP$, we get $AP=\frac{5}{\sqrt{31}}.$

 Thus the height
 \[DP = \sqrt{AD^2 - AP^2} = \sqrt{13 - \frac{25}{31}} = \sqrt{\frac{378}{31}}.\]

 
                    \begin{figure}[H]
                        \centering
                        \includegraphics[width=0.5\textwidth]{Solution/AIME_I_2023_13_Solution_Figure_2.png}
                    \end{figure}
                    

 For the other parallelepiped, using the same approach and drop a perpendicular from $E'$ onto the base at $P'$. Similarly applying Pythagorean theorem to $\triangle E'P'C'$, $\triangle E'P'A'$ and $\triangle A'P'O'$ we have
 \[C'E'^2 - C'P'^2 = A'E'^2 - A'O'^2 - (C'P'-C'O')^2.\]

 Plugging known values into the above equation and solve for the only unknown variable $C'P'$, we get $C'P'=\frac{5}{\sqrt{21}}.$

 Thus the height
 \[E'P' = \sqrt{C'E'^2 - C'P'^2} = \sqrt{13 - \frac{25}{21}} = \sqrt{\frac{248}{21}}.\]

 The ratio between the two is therefore
 \[\frac{DP}{E'P'} = \frac{\sqrt{\frac{378}{31}}}{\sqrt{\frac{248}{21}}} = \sqrt{\frac{2\cdot3^3\cdot7}{31}\cdot\frac{3\cdot7}{2^3\cdot31}} = \frac{3^2\cdot7}{2\cdot31} = \frac{63}{62}\], giving $\boxed{125}$.

 



\textbf{Solution 5 (Visual)}


                    \begin{figure}[H]
                        \centering
                        \includegraphics[width=0.5\textwidth]{Solution/AIME_I_2023_13_Solution_Figure_3.png}
                    \end{figure}
                    Let us inscribe a tetrahedron $ACB'D'$ in given parallelepiped so that its edges coincide with the diagonals of the faces of the parallelepiped. Note that the three edges outgoing from the vertex $B'$ have the same length $b$, and the three edges at the base have a different length $a.$ 

 The volume of the tetrahedron $V= \frac {a^2}{12}\sqrt{3b^2 - a^2}$ is three times less than the volume of the parallelepiped.

 In second parallelepiped $a$ and $b$ change the positions.

 Required ratio is $\frac {a^2 \cdot \sqrt {3b^2 - a^2}}{b^2 \cdot \sqrt {3a^2 - b^2}} = \frac {21 \sqrt {3 \cdot 31 - 21}}{31 \cdot \sqrt{3 \cdot 21 - 31}} = \frac {21 \sqrt {72}}{31 \sqrt {32}}= \frac {63}{62}.$

 <i>\textbf{Claim}</i>

 
                    \begin{figure}[H]
                        \centering
                        \includegraphics[width=0.5\textwidth]{Solution/AIME_I_2023_13_Solution_Figure_4.png}
                    \end{figure}
                    Let $ABCD$ be the regular pyramid, $AB = AC = AD = b, BC = BD = CD = a.$

 The area of $\triangle BCD = \frac {a^2 \sqrt {3}}{4}.$

 Height $AO^2 = AB^2 - OB^2 = {b^2 - \frac {a^2}{3}}.$

 Volume $V= \frac {a^2}{12}\sqrt{3b^2 - a^2}.$

 \textbf{vladimir.shelomovskii@gmail.com, vvsss}

 



\textbf{Solution 6}


                    \begin{figure}[H]
                        \centering
                        \includegraphics[width=0.5\textwidth]{Solution/AIME_I_2023_13_Solution_Figure_5.png}
                    \end{figure}
                    

 Let $\angle CAB$ be acute and let $X$ be the foot of the altitude from $C$ to $AB$. Notice that this also implies that $X$ is the foot of the altitude from $D$ to $AB$. Now $\sqrt{13} \cdot CX = AB \cdot CX = \frac{1}{2} \sqrt{21} \cdot \sqrt{31}$ so $CX = \frac{\sqrt{21} \cdot \sqrt{31}}{2 \sqrt{13}}$ and $DX$ is the same. $CD$ must either be $\sqrt{21}$ or $\sqrt{31}$ because it is a diagonal. If $CD = \sqrt{31}$, applying the Law of Cosines on $\angle DXC$, $\cos \angle DXC = -\frac{5}{21}$ so $\sin \angle DXC = \sqrt{\frac{416}{441}}$. If $CD = \sqrt{21}$, $\cos \angle DXC = \frac{5}{31}$, so $\sin \angle DXC = \sqrt{\frac{936}{961}}$. The ratios of the two parallelepipeds is equal to the ratios of the heights (since the bases are the same) which is equal to the ratio of the sines. Therefore it is 
 \[\frac{\sqrt{\frac{936}{961}}}{\sqrt{\frac{416}{441}}} = \sqrt{\frac{936}{416}} \cdot \frac{21}{31} = \frac{3}{2} \cdot \frac{21}{31} = \frac{63}{62}\] so the answer is $63 + 62 = \boxed{125}$.

 


	\item (\textit{AIME I 2023}) The following analog clock has two hands that can move independently of each other.

                    \begin{figure}[H]
                        \centering
                        \includegraphics[width=0.5\textwidth]{Problem/AIME_I_2023_14_Problem_Figure_1.png}
                    \end{figure}
                    
Initially, both hands point to the number $12$. The clock performs a sequence of hand movements so that on each movement, one of the two hands moves clockwise to the next number on the clock face while the other hand does not move.


 Let $N$ be the number of sequences of $144$ hand movements such that during the sequence, every possible positioning of the hands appears exactly once, and at the end of the $144$ movements, the hands have returned to their initial position. Find the remainder when $N$ is divided by $1000$.


 
 \textbf{Answer: }608



\textbf{Solution 1 (Drawing the Grid)}

This problem is, in essence, the following:  A $12\times12$ coordinate grid is placed on a (flat) torus; how many loops are there that pass through each point while only moving up or right?  In other words, Felix the frog starts his journey at $(0,0)$ in the coordinate plane.  Every second, he jumps either to the right or up, until he reaches an $x$- or $y$-coordinate of $12$.  At this point, if he tries to jump to a coordinate outside the square from $(0,0)$ to $(11,11)$, he "wraps around" and ends up at an $x$- or $y$- coordinate of $0$.  How many ways are there for Felix to jump on every grid point in this square, so that he ends at $(0,0)$?  This is consistent with the construction of the flat torus as $\mathbb Z^2/12\mathbb Z^2$ (2-dimensional modular arithmetic.  $(\mathbb{Z}_{12})^2$)

 Moving on, define a $\textit{path}$ from point $A$ to point $B$ to be a sequence of "up"s and "right"s that takes Felix from $A$ to $B$.  The $\textit{distance}$ from $A$ to $B$ is the length of the shortest path from $A$ to $B$.  At the crux of this problem is the following consideration:  The points $A_i=(i,11-i), i\in{0,...,11}$ are pairwise equidistant, each pair having distance of $12$ in both directions. 

 
                    \begin{figure}[H]
                        \centering
                        \includegraphics[width=0.5\textwidth]{Solution/AIME_I_2023_14_Solution_Figure_1.png}
                    \end{figure}
                    

 A valid complete path then joins two $A_i$'s, say $A_i$ and $A_j$.  In fact, a link between some $A_i$ and $A_j$ fully determines the rest of the cycle, as the path from $A_{i+1}$ must "hug" the path from $A_i$, to ensure that there are no gaps.  We therefore see that if $A_0$ leads to $A_k$, then $A_i$ leads to $A_{i+k}$.  Only the values of $k$ relatively prime to $12$ result in solutions, though, because otherwise $A_0$ would only lead to $\{A_i:\exists n\in \mathbb Z:i\equiv kn\quad\text{mod 12}\}$.  The number of paths from $A_0$ to $A_k$ is ${12\choose k}$, and so the answer is

 
 \[{12\choose1}+{12\choose5}+{12\choose7}+{12\choose11}=1\boxed{608}.\]

 Notes:

 - One can prove that the path from $A_{i+1}$ must "hug" the path from $A_i$ by using techniques similar to those in Solution 2.

 - One can count the paths as follows:  To get from $A_0$ to $A_i$, Felix takes $k$ rights and $12-k$ ups, which can be done in  ${12\choose k}$ ways.

 



\textbf{Solution 2 (Simple Cyclic Permutation Analysis)}

This is more of a solution sketch and lacks rigorous proof for interim steps, but illustrates some key observations that lead to a simple solution.

 Note that one can visualize this problem as walking on a $N \times N$ grid where the edges warp. Your goal is to have a single path across all nodes on the grid leading back to $(0,\ 0)$. For convenience, any grid position are presumed to be in $\mod N$.

 Note that there are exactly two ways to reach node $(i,\ j)$, namely $(i - 1,\ j)$ and $(i,\ j - 1)$. 

 As a result, if a path includes a step from $(i,\ j)$ to $(i + 1,\ j)$, there cannot be a step from $(i,\ j)$ to $(i,\ j + 1)$. However, a valid solution must reach $(i,\ j + 1)$, and the only valid step is from $(i - 1,\ j + 1)$.

 So a solution that includes a step from $(i,\ j)$ to $(i + 1,\ j)$ dictates a step from $(i - 1,\ j + 1)$ to $(i,\ j + 1)$ and by extension steps from $(i - a,\ j + a)$ to $(i - a + 1,\ j + a)$. We observe the equivalent result for steps in the orthogonal direction.

 This means that in constructing a valid solution, taking one step in fact dictates N steps, thus it's sufficient to count valid solutions with $N = a + b$ moves of going right $a$ times and $b$ times up the grid. The number of distinct solutions can be computed by permuting 2 kinds of indistinguishable objects $\binom{N}{a}$.

 Here we observe, without proof, that if $\gcd(a, b) \neq 1$, then we will return to the origin prematurely. For $N = 12$, we only want to count the number of solutions associated with $12 = 1 + 11 = 5 + 7 = 7 + 5 = 11 + 1$.

 (For those attempting a rigorous proof, note that $\gcd(a, b) = \gcd(a + b, b) = \gcd(N, b) = \gcd(N, a)$).

 The total number of solutions, noting symmetry, is thus

 
 \[2\cdot\left(\binom{12}{1} + \binom{12}{5}\right) = 1608\]

 This yields $\boxed{\textbf{608}}$ as our desired answer.

 



\textbf{Solution 3 (Matrix Analysis and Permutation)}

Define a $12 \times 12$ matrix $X$.Each entry $x_{i, j}$ denotes the number of movements the longer hand moves, given that two hands jointly make $12 \left( i - 1 \right) + \left( j - 1 \right)$ movements.Thus, the number of movements the shorter hand moves is $12 \left( i - 1 \right) + \left( j - 1 \right) - x_{i, j}$.

 Denote by $r_{i, j}$ the remainder of $x_{i, j}$ divided by 12.Denote by $R$ this remainder matrix.

 If two hands can return to their initial positions after 144 movements, then $r_{12, 12} = 0$ or 11.Denote by $S_0$ (resp. $S_{11}$) the collection of feasible sequences of movements, such that $r_{12, 12} = 0$ (resp. $r_{12, 12} = 11$).

 Define a function $f :  S_0 \rightarrow S_{11}$, such that for any $\left\{ x_{i,j} , \ \forall \ i, j \in \left\{ 1, 2, \cdots , 12 \right\} \right\} \in S_0$, the functional value of the entry indexed as $\left( i, j \right)$ is $12 \left( i - 1 \right) + \left( j - 1 \right) - x_{i, j}$.Thus, function $f$ is bijective. This implies $| S_0 | = | S_{11} |$.

 In the rest of analysis, we count $| S_0 |$.

 We make the following observations:

 
                \begin{itemize}
                    \item  $x_{1, 1} = 0$ and $12 | x_{12, 12}$. These follow from the definition of $S_0$.\item  Each column of $R$ is a permutation of $\left\{ 0, 1, \cdots , 11 \right\}$. The reasoning is as follows. Suppose there exist $i < i'$, $j$, such that $r_{i, j} = r_{i', j}$. Then this entails that the positions of two hands after the $\left( 12 \left( i' - 1 \right) + \left( j - 1 \right) \right)$th movement coincide with their positions after the  $\left( 12 \left( i - 1 \right) + \left( j - 1 \right) \right)$th movement.\item  For any $j \in \left\{ 1, 2 ,\cdots , 11 \right\}$, $x_{i, j+1} - x_{i, j}$ is equal to either 0 for all $i$ or 1 for all $i$.     The reasoning is as follows. If this does not hold and the $j$th column in $R$ is a permutation of $\left\{ 0, 1, \cdots , 12 \right\}$, then the $j+1$th column is no longer a permutation of $\left\{ 0, 1, \cdots , 12 \right\}$. This leads to the infeasibility of the movements.\item  $x_{i+1, 1} = x_{i, 12}$ for any $i \in \left\{ 1, 2, \cdots , 11 \right\}$. This follows from the conditions that the $12$th column in $R$ excluding $r_{12, 12}$ and the first column in $R$ excluding $x_{1, 1}$ are both permutations of $\left\{ 1, 2, \cdots , 11 \right\}$.
                \end{itemize}


                All observations jointly imply that $x_{i, 12} = i \cdot x_{1, 12}$.Thus, $\left\{ r_{1, 12}, r_{2, 12} , \cdots , r_{11, 12} \right\}$ is a permutation of $\left\{ 1, 2, \cdots , 11 \right\}$.Thus, $x_{1, 12}$ is relatively prime to 12.

 Because $x_{1, 1} = 0$ and $x_{1, 12} - x_{1, 1} \leq 11$, we have $x_{1, 12} = 1$, 5, 7, or 11.

 Recall that when we move from $x_{1, 1}$ to $x_{1, 12}$, there are 11 steps of movements. Each movement has $x_{1, j+1} - x_{i, j} = 0$ or 1.Thus, for each given $x_{1, 12}$, the number of feasible movements from $x_{1, 1}$ to $x_{1, 12}$ is $\binom{11}{x_{1, 12}}$.

 Therefore, the total number of feasible movement sequences in this problem is
 \begin{align*} | S_0 | + | S_{11} | & = 2 | S_0 | \\ & = 2 \cdot \sum_{x_{1, 12} = 1, 5, 7, 11} \binom{11}{x_{1, 12}} \\ & = 2 \left( 11 + 462 + 330 + 1 \right) \\ & = 1608 . \end{align*}


 Therefore, the answer is $\boxed{\textbf{(608)}}$.

 \textbf{Supplementary Contents} This can be interpreted as the 2D-model stated in Solution 1. After certain observations, we found out that the first 12 steps determines the all the rest.

 This is similar to the situation of Problem 11 of 2025 AIME II (\href{https://artofproblemsolving.com/wiki/index.php/2025\_AIME\_II\_Problems/Problem\_11}{https://artofproblemsolving.com/wiki/index.php/2025\_AIME\_II\_Problems/Problem\_11}). In this case, every possible positioning of the hands appears exactly once, meaning that the length of the repeating cycle is 12. As we can see from the pattern, only 1 and numbers relatively prime to 12 can fit in the conditions, which are 1, 5, 7, 11.

 In conclusion,
 \[{12\choose1}+{12\choose5}+{12\choose7}+{12\choose11}=1\boxed{608}.\]

 


	\item (\textit{AIME I 2023}) Find the largest prime number $p<1000$ for which there exists a complex number $z$ satisfying


 
                \begin{itemize}
                    \item the real and imaginary part of $z$ are both integers;
                \end{itemize}


                
                \begin{itemize}
                    \item $|z|=\sqrt{p},$ and
                \end{itemize}


                
                \begin{itemize}
                    \item there exists a triangle whose three side lengths are $p,$ the real part of $z^{3},$ and the imaginary part of $z^{3}.$
                \end{itemize}


                
 \textbf{Answer: }349



\textbf{Solution}

Assume that $z=a+bi$. Then,
 \[z^3=(a^3-3ab^2)+(3a^2b-b^3)i\]Note that by the Triangle Inequality,
 \[|(a^3-3ab^2)-(3a^2b-b^3)|<p\implies |a^3+b^3-3ab^2-3a^2b|<a^2+b^2\]Thus, we know
 \[|a+b||a^2+b^2-4ab|<a^2+b^2\]Without loss of generality, assume $a>b$ (as otherwise, consider $i^3\overline z=b+ai$). If $|a/b|\geq 4$, then
 \[17b^2\geq a^2+b^2>|a+b||a^2+b^2-4ab|\geq |b-4b||16b^2-16b^2+b^2|=3b^3\]`Thus, this means $b\leq\frac{17}3$ or $b\leq 5$. Also note that the roots of $x^2-4x+1$ are $2\pm\sqrt 3$, so thus if $b\geq 6$,
 \[2\sqrt 3b=(2(2-\sqrt 3)-4)b<a<4b\]Note that
 \[1000>p=a^2+b^2\geq 12b^2+b^2=13b^2\]so $b^2<81$, and $b<9$. If $b=8$, then $16\sqrt 3\leq a\leq 32$. Note that $\gcd(a,b)=1$, and $a\not\equiv b\pmod 2$, so $a=29$ or $31$. However, then $5\mid a^2+b^2$, absurd.

 If $b=7$, by similar logic, we have that $14\sqrt 3 <a< 28$, so $a=26$. However, once again, $5\mid a^2+b^2$. If $b=6$, by the same logic, $12\sqrt3<a<24$, so $a=23$, where we run into the same problem. Thus $b\leq 5$ indeed.

 If $b=5$, note that
 \[(a+5)(a^2+25-20a)<a^2+25\implies a<20\]We note that $p=5^2+18^2=349$ works. Thus, we just need to make sure that if $b\leq 4$, $a\leq 18$. But this is easy, as
 \[p>(a+b)(a^2+b^2-4ab)\geq (4+18)(4^2+18^2-4\cdot 4\cdot 18)>1000\]absurd. Thus, the answer is $\boxed{349}$.

 



\textbf{Solution 2}

Denote $z = a + i b$. Thus, $a^2 + b^2 = p$.

 Thus,
 \[z^3 = a \left( a^2 - 3 b^2 \right) + i b \left( - b^2 + 3 a^2 \right) .\]

 Because $p$, ${\textrm Re} \left( z^3 \right)$, ${\textrm Im} \left( z^3 \right)$ are three sides of a triangle, we have ${\textrm Re} \left( z^3 \right) > 0$ and ${\textrm Im} \left( z^3 \right) > 0$.Thus,
 \begin{align*} a \left( a^2 - 3 b^2 \right) & > 0 , \hspace{1cm} (1) \\ b \left( - b^2 + 3 a^2 \right) & > 0. \hspace{1cm} (2) \end{align*}


 Because $p$, ${\textrm Re} \left( z^3 \right)$, ${\textrm Im} \left( z^3 \right)$ are three sides of a triangle, we have the following triangle inequalities:
 \begin{align*} {\textrm Re} \left( z^3 \right) + {\textrm Im} \left( z^3 \right) & > p \hspace{1cm} (3) \\ p + {\textrm Re} \left( z^3 \right) & > {\textrm Im} \left( z^3 \right) \hspace{1cm} (4) \\ p + {\textrm Im} \left( z^3 \right) & > {\textrm Re} \left( z^3 \right) \hspace{1cm} (5) \end{align*}


 We notice that $| z^3 | = p^{3/2}$, and ${\textrm Re} \left( z^3 \right)$, ${\textrm Im} \left( z^3 \right)$, and $| z^3 |$ form a right triangle. Thus, ${\textrm Re} \left( z^3 \right) + {\textrm Im} \left( z^3 \right) > p^{3/2}$.Because $p > 1$, $p^{3/2} > p$.Therefore, (3) holds.

 Conditions (4) and (5) can be written in the joint form as
 \[\left| {\textrm Re} \left( z^3 \right) - {\textrm Im} \left( z^3 \right) \right| < p . \hspace{1cm} (4)\]

 
We have
 \begin{align*} {\textrm Re} \left( z^3 \right) - {\textrm Im} \left( z^3 \right) & = \left( a^3 - 3 a b^2 \right) - \left( - b^3 + 3 a^2 b \right) \\ & = \left( a + b \right) \left( a^2 - 4 ab + b^2 \right) \end{align*}
and $p = a^2 + b^2$.

 Thus, (5) can be written as
 \[\left| \left( a + b \right) \left( a^2 - 4 ab + b^2 \right)  \right| < a^2 + b^2 . \hspace{1cm} (6)\]

 Therefore, we need to jointly solve (1), (2), (6).From (1) and (2), we have either $a, b >0$, or $a, b < 0$.In (6), by symmetry, without loss of generality, we assume $a, b > 0$.

 Thus, (1) and (2) are reduced to
 \[a > \sqrt{3} b . \hspace{1cm} (7)\]

 Let $a = \lambda b$. Plugging this into (6), we get
 \begin{align*} \left| \left( \left( \lambda - 2 \right)^2 - 3 \right)  \right| < \frac{1}{b} \frac{\lambda^2 + 1}{\lambda + 1} .  \hspace{1cm} (8) \end{align*}


 Because $p= a^2 + b^2$ is a prime, $a$ and $b$ are relatively prime.

 Therefore, we can use (7), (8), $a^2 + b^2 <1000$, and $a$ and $b$ are relatively prime to solve the problem.

 To facilitate efficient search, we apply the following criteria:

 To satisfy (7) and $a^2 + b^2 < 1000$, we have $1 \leq b \leq 15$.In the outer layer, we search for $b$ in a decreasing order.In the inner layer, for each given $b$, we search for $a$.Given $b$, we search for $a$ in the range $\sqrt{3} b < a < \sqrt{1000 - b^2}$.We can prove that for $b \geq 9$, there is no feasible $a$.The proof is as follows.

 For $b \geq 9$, to satisfy $a^2 + b^2 < 1000$, we have $a \leq 30$.Thus, $\sqrt{3} < \lambda \leq \frac{30}{9}$.Thus, the R.H.S. of (8) has the following upper bound
 \begin{align*} \frac{1}{b} \frac{\lambda^2 + 1}{\lambda + 1} & < \frac{1}{b} \frac{\lambda^2 + \lambda}{\lambda + 1} \\ & = \frac{\lambda}{b} \\ & \leq \frac{\frac{30}{9}}{9} \\ & < \frac{10}{27} . \end{align*}


 Hence, to satisfy (8), a necessary condition is
 \begin{align*} \left| \left( \left( \lambda - 2 \right)^2 - 3 \right)  \right| < \frac{10}{27} . \end{align*}


 However, this cannot be satisfied for $\sqrt{3} < \lambda \leq \frac{30}{9}$.Therefore, there is no feasible solution for $b \geq 9$.Therefore, we only need to consider $b \leq 8$.

 We eliminate $a$ that is not relatively prime to $b$.

 We use the following criteria to quickly eliminate $a$ that make $a^2 + b^2$ a composite number.

 
                \begin{itemize}
                    \item For $b \equiv 1 \pmod{2}$, we eliminate $a$ satisfying $a \equiv 1 \pmod{2}$.
                \end{itemize}


                
                \begin{itemize}
                    \item For $b \equiv \pm 1 \pmod{5}$ (resp. $b \equiv \pm 2 \pmod{5}$), we eliminate $a$ satisfying $a \equiv \pm 2 \pmod{5}$ (resp. $a \equiv \pm 1 \pmod{5}$).
                \end{itemize}


                
\item For the remaining $\left( b, a \right)$, check whether (8) and the condition that $a^2 + b^2$ is prime are both satisfied.

 The first feasible solution is $b = 5$ and $a = 18$.Thus, $a^2 + b^2 = 349$.

 \item For the remaining search, given $b$, we only search for $a \geq \sqrt{349 - b^2}$.

 Following the above search criteria, we find the final answer as $b = 5$ and $a = 18$.Thus, the largest prime $p$ is $p = a^2 + b^2 = \boxed{\textbf{(349) }}$.

 



\textbf{Solution 3}

Let $z = a + b i (a,b\in I)$. $a^2 + b^2 = prime < 1000$, $(a,b) = 1$.

 According to the question, ${\textrm Re} \left( z^3 \right)$, ${\textrm Im} \left( z^3 \right)$, and $|z^3|$ construct the side-lengths of a non-degenerate triangle.

 $z^3 = (a+bi)^3 = a^3+3a^2bi-3ab^2-b^3i = (a^3 - 3ab^2) + (3a^2b-b^3)i$

 ${\textrm Re} \left( z^3 \right) = a^3 - 3ab^2 > 0 => a(a^2 - 3b^2) > 0$

 
                    \begin{figure}[H]
                        \centering
                        \includegraphics[width=0.5\textwidth]{Solution/AIME_I_2023_15_Solution_Figure_1.png}
                    \end{figure}
                    

 ${\textrm Im} \left( z^3 \right) = a^3 - 3ab^2 > 0 => a(a^2 - 3b^2) > 0$

 
                    \begin{figure}[H]
                        \centering
                        \includegraphics[width=0.5\textwidth]{Solution/AIME_I_2023_15_Solution_Figure_2.png}
                    \end{figure}
                    

 This means that the values of $a$ and$b$ should be limited in coincident areas these two graphs. 

 
                    \begin{figure}[H]
                        \centering
                        \includegraphics[width=0.5\textwidth]{Solution/AIME_I_2023_15_Solution_Figure_3.png}
                    \end{figure}
                    

 Also

 ${\textrm Re} \left( z^3 \right) + {\textrm Im} \left( z^3 \right) > |z^3| > |z||z^2|>|z^2|$

 $|{\textrm Re} \left( z^3 \right) - {\textrm Im} \left( z^3 \right)| < |z|^2$

 $|a^3-3ab^2-(3a^2b-b^3)| < a^2+b^2$

 $=|a^3-3ab^2+b^3-3a^2b|$

 $=|a^3+b^3-3ab(a+b)|$

 $=|a+b||a^2-4ab+b^2|<a^2+b^2 (*)$

 If $a<0,b>0$, $|a^2-4ab+b^2|>|a^2+b^2|$, making statement $(*)$ false.Combining with the former graph depicting possible ranges of $a,b$, by loss of generality, we assume $a,b$ both $>0$ and exists in the first $30^{\circ}$ of the circle.

 Let $\frac{a}{b} = \lambda > \sqrt{3}$.

 $(*) |b^3(1+\lambda)\cdot({\lambda}^{2}-4\lambda+1)| < b^2(1+{\lambda}^{2})$

 $b<|\frac{1+{\lambda}^2}{1+\lambda}|\cdot|\frac{1}{{\lambda}^{2}-4\lambda+1}|$

 To clearly visualize, we graph out $|\frac{1+{\lambda}^2}{1+\lambda}|$ and $|{\lambda}^{2}-4\lambda+1|$ separately.

 
                    \begin{figure}[H]
                        \centering
                        \includegraphics[width=0.5\textwidth]{Solution/AIME_I_2023_15_Solution_Figure_4.png}
                    \end{figure}
                    

 When $\lambda$ is around $2+\sqrt{3}$, b reaches its maximum upper bound. 

 $b^2(1+{\lambda}^{2}) < 1000$

 $b^2<66$

 $b\le 8$

 Testing values of $b$ in decreasing order, starting from 8, we test out each corresponding value of $a$($b\cdot\lambda$)by trying the two whole numbers closest to the real value of $a$.

 We finally get that $b=5, a=18$ and $p = 5^2+18^2 = \boxed{349}$

 


\end{enumerate}